%\documentclass[atmp]{ipart_v1}
\documentclass[reqno,11pt]{amsart}
%\documentclass[11pt]{article}

%\usepackage[notref, notcite]{showkeys}
\usepackage{graphicx}
\usepackage{amssymb}
\usepackage{amsmath}
\usepackage{amsfonts}
\usepackage{amsthm}
\usepackage{color}
\usepackage{hyperref}
\usepackage[mathscr]{eucal}
\usepackage{enumerate}
\usepackage{times}
\usepackage{algorithm}
\usepackage{algorithmic}
%\usepackage{authblk}
\usepackage[skip=10pt]{caption} % Adjust the skip value as needed


%\numberwithin{equation}{section}
\allowdisplaybreaks[1]

\newtheorem{Def}{Definition}[section]
\newtheorem{Thm}[Def]{Theorem}
\newtheorem{Prop}[Def]{Proposition}
\newtheorem{Lemma}[Def]{Lemma}
\newtheorem{Remark}[Def]{Remark}
\newtheorem{Corollary}[Def]{Corollary}
\newtheorem{Example}[Def]{Example}
\newtheorem{Assumption}[Def]{Assumption}
\newcommand{\beq}{\begin{equation}}
\newcommand{\eeq}{\end{equation}}
\newcommand{\Proof}{\begin{proof}}
\newcommand{\QED}{\end{proof} \noindent}
\newcommand{\QEDrem}{\ \hfill $\Diamond$}
\newtheorem{proposition}[Def]{Proposition}

\newcommand{\spc}{\;\;\;\;\;\;\;\;\;\;}
\newcommand{\msp}{\hspace{-.1cm}}
\newcommand{\mm}{\hspace{-.08cm}\cdot \hspace{-.08cm}}

\newcommand{\M}{\mathcal{M}}
\newcommand{\C}{\mathbb{C}}
\newcommand{\R}{\mathbb{R}}
%\newcommand{\1}{\mbox{\rm 1 \hspace{-1.05 em} 1}}
\newcommand{\Z}{\mathbb{Z}}
\newcommand{\N}{\mathcal{N}}
\newcommand{\VM}{\mathcal{V}\hspace{-.05cm}\mathcal{M}}

\newcommand{\Gammati}{\tilde{\Gamma}}
\newcommand{\Riem}{{\rm Riem}}
\newcommand{\A}{\mathcal{A}}
\newcommand{\Ati}{\tilde{\mathcal{A}}}
\newcommand{\Aop}{\mathcal{A}_\textbf{b}}
\newcommand{\X}{X}
\newcommand{\red}{\color{red}}
\newcommand{\cutoff}{\chi_{_\delta}}


\title[CayleyPy RL]{CayleyPy RL: Pathfinding and Reinforcement Learning on Cayley Graphs}

\author[A. Chervov]{A. Chervov}
\address{Institut Curie, Paris, France}
\email{al.chervov@gmail.com}

\author[A.Soibelman, ... ]{
A.Soibelman, S.Lytkin, I.Kiselev, S.Fironov, A.Lukyanenko, A.Dolgorukova, A.Ogurtsov,
F.Petrov, S.Krymskii, M.Evseev, L.Grunvald, D.Gorodkov, G.Antiufeev, G.Verbii, 
V.Zamkovoy, L.Cheldieva, I.Koltsov, A. Sychev, M.Obozov, 
A.Eliseev, S.Nikolenko, N.Narynbaev, R.Turtayev,
N. Rokotyan, S.Kovalev,
A.Rozanov, V.Nelin, S.Ermilov,
 L.Shishina, D.Mamayeva, A.Korolkova,
K.Khoruzhii, A.Romanov }

%\author[A.Chervov]{A.Chervov}
% \author[A.Soibelman]{A.Soibelman}
% \author[S.Lytkin]{S.Lytkin}
% \author[S.Fironov]{S.Fironov}
% \author[A.Lukyanenko]{A.Lukyanenko}
% \author[A.Dolgorukova]{A.Dolgorukova}
% \author[A.Ogurtsov]{A.Ogurtsov}
% \author[F.Petrov]{F.Petrov}
% \author[S.Krymskii]{S.Krymskii}
% \author[M.Evseev]{M.Evseev}
% \author[L.Grunvald]{L.Grunvald}
% \author[D.Gorodkov]{D.Gorodkov}
% \author[G.Antiufeev]{G.Antiufeev}
% \author[G.Verbii]{G.Verbii}
% \author[V.Zamkovoy]{V.Zamkovoy}
% \author[L.Cheldieva]{L.Cheldieva}
% \author[I.Koltsov]{I.Koltsov}
% \author[A.Sychev]{A.Sychev}
% \author[M.Obozov]{M.Obozov}
% \author[A.Eliseev]{A.Eliseev}
% \author[S.Nikolenko]{S.Nikolenko}
% \author[N.Narynbaev]{N.Narynbaev}
% \author[R.Turtayev]{R.Turtayev}
% \author[N.Rokotyan]{N.Rokotyan}
% \author[S.Kovalev]{S.Kovalev}
% \author[A.Rozanov]{A.Rozanov}
% \author[V.Nelin]{V.Nelin}
% \author[S.Ermilov]{S.Ermilov}
% \author[I.Kiselev]{I.Kiselev}
% \author[L.Shishina]{L.Shishina}
% \author[D.Mamayeva]{D.Mamayeva}
% \author[A.Korolkova]{A.Korolkova}
% \author[K.Khoruzhii]{K.Khoruzhii}
% \author[A.Romanov]{A.Romanov}


% \author[B.\ Temple]{Blake Temple \\ \\ January 28, 2024}
% \address{Department of Mathematics\\ University of California\\ Davis, CA 95616\\ USA}
% \email{temple@math.ucdavis.edu}





\includeonly{}

\begin{document} 

\begin{abstract}
This paper is the second in a series of studies on developing efficient artificial intelligence-based approaches to pathfinding on extremely large graphs (e.g., $10^{70}$ nodes) focusing on Cayley graphs and mathematical applications. The open-source CayleyPy project is a central component of our research. The present paper proposes a novel combination of a reinforcement learning approach with a more direct diffusion distance approach from the first paper. Our analysis includes benchmarking various choices for the key building blocks of the approach: neural network architectures, random walk generators, and beam search pathfinding. We compared these methods against the classical computer algebra system GAP, demonstrating that they "overcome the GAP" for the considered examples.
We examine the Cayley graph of the symmetric group with cyclic shift and transposition generators as a particular mathematical application. We strongly support the OEIS-A186783 conjecture that the diameter equals n(n-1)/2 by machine learning and mathematical methods. We identify the conjectured longest element and generate its decomposition of the desired length. We prove a diameter lower bound of $n(n-1)/2-n/2-1$ and an upper bound of $n(n-1)/2+3n$ by presenting the algorithm with given complexity. We also present several conjectures motivated by numerical experiments, including observations on the central limit phenomenon (with growth approximated by the Gumbel distribution), the uniform distribution for the spectrum of the graph, and a numerical study of sorting networks.
We create challenges on the Kaggle platform to stimulate crowdsourcing activity and invite contributions to improve and benchmark approaches to Cayley graph pathfinding and other tasks.
\end{abstract}



\keywords{Machine learning, reinforcement learning, Cayley graphs}

\maketitle 

\setcounter{tocdepth}{1}
\small
\tableofcontents
\normalsize

\section{Introduction}

\subsection{Broader context}
Deep learning has revolutionized various fields of research and was recognized with several Nobel Prizes in 2024.  
In recent years, machine learning has been emerging as "a tool in theoretical science"~\cite{douglas2022machine}, leading to several noteworthy applications to mathematical problems:~\cite{lample2019deep,davies2021advancing, bao2021polytopes, 
romera2024mathematical,
coates2024machine,alfarano2024global, charton2024patternboost,shehper2024makes,swirszcz2025advancing}.  %{hashemi2025transformers}.

The paper presents advancements in applying artificial intelligence methods to the mathematical problem of Cayley graph pathfinding, equivalent to decomposing a group element into a product of generators. It is the second work in a series devoted to the CayleyPy project to develop efficient machine learning-based approaches and an open-source library to deal with extra-large graphs.
(e.g. $10^{70}$) with a focus on finite Cayley graphs and mathematical applications.

From a broader perspective, pathfinding is a specific case of the planning problem, where one must determine a sequence of actions to transition between two given states. Similar challenges arise in robotics and games like chess or Go, where the objective is to plan moves that lead from the initial position to a winning position. Mathematically, such problems are modeled as pathfinding on graphs (state transition graphs), where nodes represent possible states, and edges correspond to state transitions based on actions or moves. The planning task then reduces to finding a path from a given initial node to one or more target nodes. The breakthrough works AlphaGo and AlphaZero~\cite{silver2016mastering}, 
\cite{silver2017mastering} by Google DeepMind have demonstrated that machine learning can significantly outperform previously known methods. They serve as both a prototype and an inspiration for many developments, including ours. The method consists of two main components: a neural network trained to suggest moves from any given state and a graph pathfinding algorithm that helps to find a path even when the neural network's predictions are not sufficiently precise. The core idea is similar across most approaches, but achieving optimal performance for each task requires analyzing and benchmarking various neural networks and graph pathfinding methods. 

\begin{table}[ht!]
    \centering
    \begin{tabular}{|c|c|c|}
        \hline
        Graph  & Group  & RL  \\
        \hline
        vertex / node & permutation / element & state \\
        edge & generator & action \\
        path & word  & trajectory \\
        shortest path & shortest word  & maximal return \\
        diameter & longest element & ?\\
        \hline
    \end{tabular}
    \captionsetup{skip=10pt} 
    \caption{Comparison of graph theory, algebraic and reinforcement learning terminology.}
    \label{tab:terms}
\end{table}

One of the goals of this work is to apply these methods to a particular class of Cayley graphs. As in our previous paper~\cite{chervov2025machinelearningapproachbeats}, we choose a set of \textit{generators} $a_1^{\pm 1},\ldots, a_k^{\pm 1}\in S_n$. All finite products (words) of these elements form a subgroup $G \subset S_n$ which we identify with its Cayley graph. Due to vertex-transitivity of Cayley graphs, the pathfinding problem can be reduced to searching a path from some element $g\in G$ to the identity permutation $e$; this is the same as finding a decomposition of $g$ into a word $\prod\limits_{j=1}^m a_{i_j}^{\pm 1}$, $1 \leq i_j \leq k$. The problem of finding the shortest path (or word) can also be reformulated in terms of a generic reinforcement learning approach discussed in Section~\ref{Bellman_section}. A set of dictionary between RL, graphs, and group  is presented in Table~\ref{tab:terms}.
% A set of interchangeable terms related to the paths in the Cayley graphs of symmetric subgroups is presented in Table~\ref{tab:terms}.




\subsection{ Present case study - LRX Cayley graph.   } % Approaches and benchmarks

\begin{figure}[ht!]
 \centering
 \includegraphics[width=1.0\linewidth]{figs/LRX.jpg}
 \caption{LRX Cayley graph for $S_5$}\label{fig:LRX1}
\end{figure}

The paper presents both AI contributions and mathematical results and conjectures. We use \textit{LRX Cayley graphs} as primary benchmark in the subsequent discussion.

{\bf LRX Cayley graphs.} 
%Here, we work with a particular Cayley graph defined for any symmetric group $S_n$. 
For each $n\in\mathbb N$ the symmetric group $S_n$ can be generated by three permutations:

\begin{enumerate}
    \item 
    left cyclic shift $L = (1\ 2\ 3 \ \ldots\ n-1 \ 0)$;
    \item
    right cyclic shift $R = L^{-1} = (n-1\ 0\ 1 \ 2\  \ldots\ n-3 \ n-2)$;
    \item 
    transposition (eXchange) of the first two positions $X = X^{-1} = (0,1)$.
\end{enumerate}
%the two cyclic shifts -- left (L) and right (R) -- and exchange (X) of the first two positions. 
These generators and the corresponding Cayley graph will be denoted LRX following \href{https://oeis.org/A186783}{OEIS-A186783}.
%The image of LRX Cayley graph for $S_5$ is shown in Fig.~\ref{fig:LRX1}.
Visualization (Fig.~\ref{fig:LRX1}) of the LRX Cayley graph for $S_5$ reveals its complicated structure. 
%More visualizations by the minimal distortion method can be found in this \href{https://www.kaggle.com/code/iggisv9t/drawing-cayley-graphs-with-mds}{notebook}, and by originally modified force directed layout \href{https://projects.interacta.io/graphs/IM_S_N_4.csv}{here}.
LRX Cayley graph can be considered as a relative of the \href{https://www.sfu.ca/~jtmulhol/math302/puzzles-ot.html}{\textit{Oval Track}} (or \textit{TopSpin}) and \href{https://www.sfu.ca/~jtmulhol/math302/puzzles-hr.html}{\textit{Hungarian Rings}} puzzles (Fig.~\ref{fig:LRX_TopSpin_HR}).  
%Despite its simple structure, it seems to be not fully understood yet. In particular, according to \href{https://oeis.org/A186783}{OEIS-A186783} the diameter conjecture $n(n-1)/2$ remains open. 

\begin{figure}[ht!]
   \centering
   \includegraphics[width=1.0\linewidth]{figs/LRX_TopSpin_HR.jpg}
   \caption{LRX and related mechanical puzzles}\label{fig:LRX_TopSpin_HR}
\end{figure}

% Here, we work with a particular Cayley graph defined by the L - left cyclic shift, R - right cyclic shift and X - eXchange of positions $0,1$, it will be denoted LRX 

{\bf Conjectures on diameters of LRX and LX (V.M. Glushkov's 1968 problem).} 
It is conjectured \href{https://oeis.org/A186783}{OEIS-A186783} that diameters of LRX are $n(n-1)/2$,
that conjecture will be in the focus of the present paper. 
The case of two generators just and L and X (not inverse closed) goes back at least to 1968 influential paper by
``\href{https://en.wikipedia.org/wiki/Victor_Glushkov}{one of the founding fathers of Soviet cybernetics}'' V.\,M.~Glushkov 
\cite{glushkov1968completeness}, which has been much studied (survey: \cite{glukhovzubov1999lengths} pages 18-21),
but not resolved. The conjecture on the diameter in that case is presented in our subsequent paper: \cite{CayleyPyGrowth}.

Not exactly in the present scope, but worth mentioning: there was a long-standing problem resolved positively recently \cite{SawadaWilliams2019SigmaTau} of whether the directed graph generated by L and X (also denoted Sigma and Tau)
has a hamiltonian path.  D.Knuth assigned to that problem complexity 48 out of 50, it was first articulated in
\cite{NijenhuisWilf1975CombinatorialAlgorithms}.

\subsection{ AI contributions.  }~ % Approaches and benchmarks

{\bf Development of AI methodology.} We suggest a modified deep Q-learning (DQN) method for the case of the graph pathfinding task, which resolves the so-called sparse reward problem of the standard DQN approach for graph pathfinding. The proposal combines the DQN and the diffusion distance approach from~\cite{chervov2025machinelearningapproachbeats} and presents an analysis and benchmarks of these techniques, including tests of various neural network architectures such as MLP, CNN, and transformers. %An analysis and benchmarks of these approaches are presented, including tests of various neural network architectures such as MLP, CNN, and transformers. 
These methods have been shown to find paths for LRX Cayley graphs up to $n$ around 40 and with small modification (X-trick, which is specfic to LRX) even beyond 100.
%- for the particular class of generators (LRX), which are the main focus of this paper. 
We demonstrate that they "overcome the GAP": outperform the classical computer algebra system GAP on the same task, which can work only up to $n$ around 20, also producing much more optimal paths and working faster.

{\bf Efficient PyTorch implementation.} We devote significant attention to developing an original, highly optimized code. The current version is specific to Cayley graphs of permutation groups (matrix groups will be supported later). We propose many technical solutions for fast parallel processing of large amounts of permutations, including product computations, efficient hashing, and extracting unique states. The code works on both CPU and GPU without modification.

%{\bf One line of code bringing orders of magnitudes improvement.} 
{\bf X-trick.} We report a curious finding: a single-line modification (which we call \textit{X-trick})  in the code of our main graph pathfinding module (beam search) dramatically improves performance. This enhancement increases the feasible pathfinding size
%from symmetric groups $S_n$ with $n$
to $n$ around 100 (apparently even further with more powerful hardware). Unfortunately, this improvement is currently specific to the particular LRX Cayley graphs considered in this paper, but we hope it can be generalized. 
Moreover, we provide a mathematical counterpart of that finding - X-trick modified random walks are mixing much faster - see discussion below. 

\subsection{Mathematical contributions.} 

We present several conjectures about LRX Cayley graphs which could leverage a more general understanding of the properties of that Cayley graph and its relatives.


{\bf Diameter conjecture and the longest element.} Surprisingly, according to \href{https://oeis.org/A186783}{OEIS-A186783}, it is an open conjecture that the diameter of the Cayley graph is $n(n-1)/2$ for $S_n$. 
%These generators and Cayley graphs will be denoted LRX as in OEIS.
%{\bf The longest element.} 
By implementing an efficient brute-force breadth-first search algorithm, we traverse the full LRX graphs for $n\le 15$, confirming that the longest element has a length of $n(n-1)/2$, as predicted. Moreover, we observe that this element $\ell_n$ is always unique and follows a clear pattern

\begin{equation}
\label{eq:longest}
    \ell_n = (1 \ 0 \ n-1\ n-2\ \ldots\ 3\ 2) = (0,1)(n-1,2)(n-2,3)\ldots = R^2 e^{\leftarrow},
\end{equation}
where $e^{\leftarrow} = (n-1 \ n-2 \ \ldots \ 2 \ 1\ 0)$ is reversed identity permutation.
%--- it is a product of transpositions: $(0,1)(n-1,2)(n-2,3)(n-3,4)\ldots$, 
We expect~\eqref{eq:longest} holds for all $n$. Interestingly, the longest element $\ell_n$ is equal not to the reversed identity permutation itself but to the element obtained by two right shifts from it.

{\bf AI and crowd-sourcing to support the conjecture.} For the longest element $\ell_n$, we have launched various versions of our pipelines, and have consistently found that a decomposition is never shorter than $n(n-1)/2$. We organized a \href{https://www.kaggle.com/competitions/lrx-oeis-a-186783-brainstorm-math-conjecture/overview}{challenge} on the Kaggle platform suggesting to decompose $\ell_n$, $n=2, \ldots, 200$, into the shortest possible word.
%giving that particular element for $n=2, \ldots, 200$ and requesting decompositions of the shortest possible length. 
Once again, no participant found a decomposition shorter than $n(n-1)/2$, thereby confirming the expectation.

{\bf Rigorous lower and upper bounds for the diameter.} We prove the lower bound $n(n-1)/2 - n/2-1$ by a combinatorial argument (improving $n(n-1)/3$~\cite{babai1989small}); later, we prove the upper bound $n(n-1)/2 +3n$ by the algorithm with such complexity (improving $\frac 32 n^2$~\cite{kuppili2020upper}). %To the best of our knowledge, these are the best estimates and algorithm. 
We also develop another algorithm that empirically shows better complexity $n(n-1)/2 + n/2$, but it lacks a rigorous proof for such an estimate.  
Finally, we present an explicit decomposition for the conjecturally longest element $\ell_n$ of the desired length $n(n-1)/2$. 

{\bf Central limit phenomena for growth.} 
Computations suggest the following conjecture: for large $n$, the growth function of LRX graphs follows a left-skewed \href{https://en.wikipedia.org/wiki/Gumbel_distribution}{Gumbel distribution}. This statement is in the vein of field-shaping works by P. Diaconis et. al.~\cite{diaconis1977spearman, diaconis1988metrics,chatterjee2017central}, who demonstrated that the growth of the neighbor transposition graph and many related statistics asymptotically follows the Gaussian normal distribution.
However we observe that skewness on the growth distribution for LRX does not tend to zero, thus Gaussian is not a choice, while Gumbel provides a good numerical fit. We also present a similar analysis for the Schreier coset graph for the action of the LRX generators on binary strings consisting of $n/2$ zeros and ones (for significantly larger range of $n$).

{\bf Random walks and mixing time.}
P.Diaconis and L.Saloff-Coste~\cite{diaconis1993comparison} obtained an important result that mixing time for LRX
is between $O(n^3)$ and $O(n^3log(n))$. We support their result numerically and complement it by considering two other types
of random walks: non-backtracking and X-trick modified. Both show faster mixing, however non-backtracking is only faster 
by a constant  factor, while X-trick leads to around $O(n^2)$ mixing. Which particially explains its effectivity in beam search.
% Numerical analysis of random walks shows that the mixing time for the LRX Cayley graph on $S_n$ is $O(n^3)$. 


{\bf Spectrum and sorting networks.}
We study numerically the spectrum of the LRX graph and observe a surprisingly uniform distribution of eigenvalues. Investigating statistics of the paths from $\ell_n$ to $e$ ("sorting networks"), we observe numerically that the pattern differs from the famous "sine curves" for Cayley graph of neighbor transpositions~\cite{angel2007random}.


%We also present a similar analysis for the Schreier coset graph for the action of the LRX generators on binary strings consisting of $n/2$ zeros and ones (for even $n$). %(Might be thought as an analogue of Grassmanian $Gr(n,2n)$ over the field with one element.)

{\bf In summary}, our analysis demonstrates that computational experiments, especially those involving AI, can be very helpful in advancing mathematical research on Cayley graphs. AI methods significantly outperform classical methods of computer algebra systems such as GAP, performing well for $S_n$ with $n$ around 40 and with a small addition of prior knowledge up to 90, while classical methods perform well only up to 20. 

\subsection{Review. Finite Cayley graphs: pathfinding techniques, results and open problems}

Let us briefly summarize existing techniques for Cayley graphs pathfinding, and provide more general context on Cayley graphs research.

{\bf Optimal pathfinding is typically NP-hard.} Finding the shortest paths on generic finite Cayley graphs is an NP-hard problem~\cite{even1981minimum} (even P-space complete \cite{jerrum1985complexity}). As is the case for many specific group families, such as $N\times N \times N$ Rubik's Cube groups~\cite{demaine2017solving} and others~\cite{bulteau2015pancake}.
It is hard to develop a practical optimal algorithm for particular groups, for instance, the first optimal algorithm for the standard $3\times 3 \times 3$ Rubik's cube has been proposed only in 1997~\cite{korf1997finding}. It required 82 Megabytes of precomputed data and 12-24 hours to solve a single cube (later improved~\footnote{M.Reid: https://www.cflmath.com/Rubik/optimal\_solver.html, H.Kociemba 2021: https://pypi.org/project/RubikOptimal/, Andrew Skalski:  https://github.com/Voltara/vcubem , \cite{Rokicki2014Diameter} }%\cite{Reid1997Optimal,HKociemba2021,Rokicki2014Diameter,Skalskivcube} 
to several cubes per second). Developing a practical optimal solvers for a $4\times 4 \times 4$ Rubik's cube is challenging open problem. 

{\bf Non-optimal pathfinding: Schreier-Sims algorithm, GAP implementation, non-effectivity for large groups.} 
The Schreier-Sims algorithm~\cite{sims1970computational} (and its randomized version \cite{knuth1991efficient}) is the method typically used for decomposing group elements into products of generators for permutation groups. It is implemented in the GAP computer algebra system. However, its outputs "are usually exponentially long"~\cite{fiat1989planning}, and that forbids practical computations for large groups (typically the sizes beyond $10^{30}$--$10^{40}$ are out of reach -- examples and analysis are given below). 


{\bf Non-optimal pathfinding: "small support" Kaggle Santa 2023 methods (generality issue).} The methods which can deal with extremely large groups (e.g., $>10^{1000}$ in size like $33\times33\times33$ Rubik's cube) were demonstrated by top participants of the recent Kaggle challenge \href{https://www.kaggle.com/competitions/santa-2023}{Santa 2023}. However, the generality of that approach is unclear.
%What is not clear whether  such methods can be applied only to some specific choices of generators or extended to any choices. 
That is in contrast to AI methodology, which is, from the beginning general not even restricted to Cayley graphs. The idea of "Santa methods" is quite noteworthy: one should first find "small support elements" achievable from original generators and take them as new generators, then one simply runs beam search with Hamming distance heuristics to find path. This method becomes effective since for small support generators, Hamming distance, and true graph distance is quite related for obvious reasons. "Small support" is those permutations that change only a small number of positions. They are well known for Rubik's cube solvers -- "commutator" or "Pif-Paf" moves~\cite{mulholland2016permutation}, \href{https://math.stackexchange.com/q/4962862/21498}{M.Brandenburg post}. In mathematics, small support elements play an important role in deep works by Babai, Seress, Helfgott, et.al.~\cite{babai1988diameter, babai2004diameter, bamberg2014bounds, helfgott2014diameter, helfgott2019growth} (slides: \href{https://www.math.auckland.ac.nz/~conder/SODO-2012/Seress-SODO2012.pdf}{A.Seress},\href{https://simons.berkeley.edu/sites/default/files/docs/6206/symtalk.pdf}{H.Helfgott}). However, the question whether one can effectively find small support elements from given generators remains open, thus making the generality of "small support methods" unclear.

{\bf Non-optimal pathfinding: algorithms for specific families of generators.} 
Some algorithmic solvers can effectively find non-optimal solutions for some particular generator choices. It seems unlikely that one can produce similar algorithms for any generators. The prototypical example is the bubble sort algorithm, which solves the problem for the neighbor transposition generators of $S_n$. Algorithmic solvers exist for Rubik's cube of arbitrary sizes and other puzzle-related groups. The other families of generators  and algorithms are important in bioinformatics for estimating the evolutionary distance. These generators are related to flips of subsequences (see, e.g.,~\cite{Pevzner1995human2mice, Pevzner1999cabbage2turnip, bulteau2019parameterized}). A related example is \href{https://en.wikipedia.org/wiki/Pancake_sorting}{pancake sorting} (or prefix sorting) problem with the first algorithm proposed by Bill Gates et.al. \cite{gates1979bounds}.

{\bf AI-methods for pathfinding.} 
To the best of our knowledge, there is no systematic effort for Cayley graph pathfinding with AI methods except for our CayleyPy project. However, for some particular cases, for example, Rubik's cube, there are several approaches, e.g., based on genetic algorithms:~\cite{swita2023solving}.%, which can be modified for the general case. 
The most notable are: DeepCube~\cite{mcaleer2019solving, agostinelli2019solving, khandelwal2024towards,agostinelli2024q} --- the first AI based successful solver for $3\times 3\times 3$ cube (\href{https://deepcube.igb.uci.edu/}{website}), the second one: EfficientCube~\cite{takano2023selfsupervision}, some other:~\cite{brunetto2017deep,johnson2021solving,amrutha2022deep,noever2022puzzle,chasmai2022cubetr,bedaywi2023solving,pan2021fourier} which not achieve complete solution. Noteworthy idea~\cite{pan2021fourier} is to combine neural networks with the representation theory of the symmetric group ---  neural net  predicts the coefficients of the non-abelian Fourier transform for the distance function. The rationale is: observed sparsity (bandlimitedness) of Fourier transform of the common distance functions on $S_n$ ~\cite{swan2017harmonic}. The recent seminal achievement from S.Gukov's team: ~\cite{shehper2024makes} creates an effective AI-based pathfinding approach for an infinite group of Andrews-Curtis moves and resolves the  Akbulut-Kirby conjecture remained open for 39 years. AI methods are also used for pathfinding in the context of planning movements in obstacle-rich environments and road networks~ \cite{pandy2022learning, kirilenko2023transpath}.

After the first version of the present work, several new contributions appeared:   \cite{douglas2025diffusion}  proposed a novel diffussion-like approach to that problem, considered several classes of permutation and matrix groups and also demonstrated that the deep learning approach works successfully not only on Rubik's cube, but on general Cayley graphs up to certain sizes.
\cite{ZiarkoBortkiewiczZawalskiEysenbachMilos2025CRTR} provided solution of 333 Rubik's cube based on 
approach of latent representations. 

{\bf Cayley graphs: applications, diameters, random walks, open conjectures.}
Let us provide more general context on Cayley graph research.
Cayley graphs are fundamental in group theory~\cite{gromov1993geometric},\cite{tao2015expansion}, and have various applications: bioinformatics
~\cite{Pevzner1995human2mice, Pevzner1999cabbage2turnip,  bulteau2019parameterized}; processors interconnection networks~\cite{akers1989group, cooperman1991applications,heydemann1997cayley}; in coding theory  and cryptography~\cite{hoory2006expander,zemor1994hash,petit2013rubik}; in quantum computing~\cite{ruiz2024quantum,sarkar2024quantum,dinur2023good, acevedo2006exploring,gromada2022some}, etc.

Pure mathematical research on finite Cayley graphs focuses on diameter estimation and the behavior of random walks. In the case of the symmetric group $S_n$, two open conjectures are easy to formulate, but somewhat representative of the field: 

\begin{itemize}
    \item 
    {\bf Babai-like conjecture:} for any choices of generators diameter of  $S_n$ is $O(n^{2})$ (see, e.g., ~\cite{helfgott2014diameter},\cite{helfgott2019growth},\cite{helfgott2015random}); 

    \item 
    {\bf Diaconis ~\cite{diaconis2013some} conjecture:}  the mixing time for random walks is $O(n^{3}\log n)$ (again for any choices of generators).
\end{itemize}




The first one can be thought of as a particular case of the Babai conjecture~\cite{babai1992diameter}, \cite{tao2015expansion}, which predicts that diameter of simple groups is not as large as one may expect: $O(\log^c|G|)$ for some absolute $c$ --- in contrast to abelian groups where it can be $O(|G|)$. This conjecture represents an intriguing connection between an algebraic property of a group — simplicity, which is independent of the choice of generators — and a property of its Cayley graphs — having small diameter, which does depend on the generators.

%This, conjecture represents an intriguing connection between algebraic property of a group (which does not depend on the choice of generators) - being simple and the property of Cayley graphs - having short diameter (which is generator dependent property). 



%Cayley graphs are fundamental for geometric group theory~\cite{gromov1993geometric}, which mainly focuses on infinite groups, where some properties of graphs (like growth) do not depend on choices of generators and represent internal properties of groups.

%Cayley graphs of finite groups are also active field of research in pure mathematics~\cite{tao2015expansion}, and have various applications:in  bioinformatics for estimating evolutionary distance~\cite{Pevzner1995human2mice},\cite{Pevzner1999cabbage2turnip},\cite{wilson2024cayley},\cite{bulteau2019parameterized}; processors interconnection networks~\cite{akers1989group},\cite{cooperman1991applications}, \cite{heydemann1997cayley};in coding theory for construction of expander graphs and related codes~\cite{hoory2006expander};in cryptography to construct specific hash functions~\cite{zemor1994hash},\cite{petit2013rubik};in machine learning~\cite{wilson2024cayley}; in quantum computing~\cite{ruiz2024quantum},~\cite{sarkar2024quantum},~\cite{dinur2023good},~\cite{acevedo2006exploring},~\cite{gromada2022some}.


Estimating diameters of Cayley graphs is a hard problem. One of its applications is latency estimation for the processor interconnection problem~\cite{akers1989group},\cite{cooperman1991applications}, \cite{heydemann1997cayley} since diameter corresponds to maximal latency. 
It is also a question of interest for community of puzzles enthusiasts --- like Rubik's cube. In the context of puzzles the diameter is often called the "God's number" --- the number of moves the best algorithm can solve the worst configuration. It is widely unknown for most of the puzzles. Its determination for the Rubik's cube $3\times 3 \times3$  for the standard choices of generators required significant efforts and computational resources~\cite{Rokicki2014Diameter}.  It is not known precisely for higher cubes, neither for $3\times 3 \times3$ Rubik's with \href{https://www.speedsolving.com/wiki/index.php?title=Metric#STM}{STM, ATM} metrics.
% cube with a little modified definitions of generators.  
Precise conjectures on $4\times 4 \times4$  and $5\times 5 \times5$ Rubik's cube diameters appeared recently~\cite{hirata2024probabilistic}, collections of conjectures and estimates for many puzzles can be found on \href{https://www.speedsolving.com/wiki/index.php?title=God%27s_Algorithm#Table_of_God.27s_Numbers}{speedsolving website}. 

As we will discuss below, novel AI methods might provide insights into such questions through computational experiments.

\section{Reinforcement learning (RL) and graphs pathfinding}
The last subsection (\ref{sec-RLDD}) presents our new proposal for combining reinforcement learning (RL) training with pretraining based on the diffusion distance approach introduced in our previous work \cite{chervov2025machinelearningapproachbeats}, which is briefly reviewed in Section \ref{sec-DD}. Both methods rely on beam search, recalled in Section \ref{sec-Beam}. The first subsection (\ref{Bellman_section}) provides a brief overview of reinforcement learning and explains how it naturally fits into the graph pathfinding framework — a simple yet not thoroughly discussed 
%surprisingly underexplored
perspective in the existing literature, to the best of our knowledge.

% Last subsection \ref{sec-RLDD} is devoted to our new proposal on combining RL-training with pretraining based on the diffusion distance approach from our previous paper \cite{chervov2025machinelearningapproachbeats}, which is reminded in  \ref{sec-DD}. Both approaches
% use beam-search - reminded in \ref{sec-Beam}. First subsection \ref{Bellman_section} is devoted to reinforcement learning reminder
% and how 
% it naturally fits into graph pathfinding setup, which is quite simple, but not thoroughly discussed in the literature to the best of our knowledge.


%\subsection{Reminders: RL, sparse reward problem, Bellman equation for graph distances,  Q-algorithm, pathfinding}

\subsection{RL setup for pathfinding and Bellman equation}
\label{Bellman_section}

%{\bf RL in brief.}
Reinforcement learning is an important subfield
of AI, which is different from supervised and unsupervised learning in the following respects: it aims to optimize {\bf cumulative rewards}, and it requires the agent to, in a sense, "{\bf create its own training set}". The agent should choose the best strategies to "explore the space" and select training examples by himself, in contrast to supervised learning, where the training set is given a priori.
Here we will briefly discuss it from the graph theory perspective.

{\bf Bellman equation and shortest paths}. Let $d(s)$ be the \textit{distance function} — the length of the shortest path from the state $s$ to the target node $e$. Denote the set of all neighbors of $s$ as $N(s)$. 
%From the subpaths optimality principle \cite[Lemma 24.1]{cormen09} we observe that if $(e, s_1, \ldots, s_{k-1}, t, s)$ is a shortest path from $e$ to $s$, then so is the path $(e, s_1, \ldots, s_{k-1}, t)$ from $e$ to $t$. Hence, the distance function $d(s)$ satisfies the \textit{Bellman equation}
It is obvious that, the distance function $d(s)$ satisfies the \textit{Bellman equation}: 

\begin{equation}
\label{bel_eq}
d(s) = 1 + \min\limits_{t \in N(s)} d(t)  
\end{equation}
and the \textit{boundary condition} 

\begin{equation}
\label{boundary_cond}
d(e) = 0.
\end{equation}

The equation is elementary but plays an important role in reinforcement learning and its deep learning generalizations discussed below.

The Bellman equation implies a simple proposition, which is easy to see inductively considering neighbor nodes to $e$, then neighbors of neighbors, etc. (actually, this is how the breadth-first search algorithm finds the shortest path).

\begin{proposition}
    For any finite connected graph $G$, the only solution to the equations~\eqref{bel_eq} and~\eqref{boundary_cond} is given by the distance function $d(s)$. Moreover, the only solutions to Bellman's equation~\eqref{bel_eq} (without boundary condition) correspond to the distance to some sets of nodes in $G$ (up to adding a constant). That set can be recovered as minimums of $d(s)$. 
\end{proposition}


{\bf RL "equivalence" to graph pathfinding.} 
Here we remind the standard form of Bellman equation for the \textit{value function} $v_*(s)$  (\cite[p.62]{sutton2018})  and relate it to the form given above.  Taking any action $a$ in state $s$ the agent receives a reward (penalty) $r(s, a)$,  and denote \textit{ discount factor} to be  $\gamma \in(0,1]$, then   \textit{Bellman optimality equation}:

% The problem of finding the shortest path in a Cayley graph can be naturally stated as a reinforcement learning problem (see also Table~\ref{tab:terms}). After taking any action $a$ in state $s$ the agent receives a reward (penalty) $r(s, a) = -1$; the only exception is the terminal (target) state $e$ with zero reward. In such setting, to maximize the cumulative return is obviously the same as to find the shortest path in the graph. To justify this, introduce the \textit{value function} $v_*(s)$ (\cite[p.62]{sutton2018}) which shows how good the state $s$ is. This function satisfies the \textit{Bellman optimality equation}

$$
 v_*(s) = \max\limits_a \big\{r(s, a) + \gamma \mathbb E_{s'} v_*(s')\big\}
 $$

\noindent
 % where $\gamma \in(0,1]$ is the \textit{ discount factor}. Our environment is deterministic, hence, $\mathbb E_{s'} v_*(s') = v_*(a\circ s)$ and
Considering environment to be deterministic,  $\mathbb E_{s'} v_*(s') = v_*(a\circ s)$ and taking $r(s, a)  = -1$, one gets:

\begin{equation}
 \label{eq:v-star-bellman}
 v_*(s) =
 -1+ \gamma\max\limits_a \big\{v_*(a\circ s)\big\}.
 \end{equation}

\noindent
Now observe that~\eqref{bel_eq} turns into~\eqref{eq:v-star-bellman} if $\gamma = 1$ and $d(s) = -v_*(s)$. By definition of the value function $v_*(e) = 0$ which is the same as~\eqref{boundary_cond}.



%The setup of reinforcement learning is very similar to the setup of graph pathfinding: states are nodes of the graph, actions correspond to edges/transitions to neighbor states by these edges, and rewards correspond to weights on the edges. The cumulative reward corresponds to length of paths (with respect to edges weights). Typically, in graph theory, shortest paths are of interest, which means that the reward is minimized, not maximized as in RL, so it is better to use term penalty instead of reward. 

% {\bf Unweighted graphs and sparse reward problem.}
% The sparse reward problem is famous in RL, and it exactly appears in Cayley graphs pathfinding and more generally for pathfinding on any unweighted graph, thus making standard RL approaches poorly performing.
% The sparse reward problem is a situation where the reward is given very rarely; thus, generating the data for training is not easy. That situation exactly happens for graph pathfinding on large graphs - where one needs to find a path to, say, one selected node out of trillions of other nodes. The reward is given only by finding the desired node - but it is practically impossible to find it with a "cold start" randomly. Pay attention that condition on all weights equal to one - play a role - if weights would be different - one may hope the strategy to choose the step with the best local reward is reasonable; however, when all weights are precisely the same, there is no local preference to make a right move.
% Several strategies exist to resolve that issue, and we propose a new one below.

% {\bf Bellman value equation  is equation for distance functions.}
% The value function (or "position estimation") in RL exactly corresponds to the distance functions on graphs. 
% To make that analogy more precise: consider an unweighted  graph $G$ (all edges weights are equal to 1), choose some node $e$ (e.g. identity of a group in Cayley graph), consider a function $d(g)$ --- the distance of the shortest path from $g$ to $e$ (it is "\href{https://en.wikipedia.org/wiki/Word_metric}{word metric}" in group theory). Then there is a simple relation:

% \begin{align} 
% %$$ 
% \label{bel_eq}
% d(g) & = 1 + \min_{n:~neigbours~of~g} d(n)  \text{~~~~~~~~~ Bellman equation}\\
% %$$
% d(e) & = 0 \text{~~~boundary condition}
% \end{align} 



% {\bf Sense of Bellman equation - local condition for globally (cumulatively) defined object.}
% The Bellman equation is an equation for the length of the shortest path --that is, it is a globally/cumulatively defined object, but the condition itself is local -- only at the neighbor of each point.  

%{\bf Idea of Q-learning and deep Q-learning algorithms for value/distance function.}
%These algorithms provide a way to compute solutions of the Bellman equation, which means to find the distance functions on graphs. One starts with some  initialization for $d(g)$, then just  the new approximation $d_k(g) = 1+\min_{n:~neigbours~of~g} d_{k-1}(n), d_{k}(e) =0$ by applying the right hand side of the equations \ref{bel_eq}. There are theoretical guarantees that such a process will converge to the solution of the Bellman equation (and boundary condition). 

%\subsection{Reinforcement learning on "small graphs"}
 

{\bf Solving Bellman - simple iteration method (dynamic programming).} 
Bellman equation on unknown function $d$ has the form $d = F(d)$, thus to solve it one can use method of simple iteration:
$d_{k+1} = F (d_{k} )$. 
Thus, solving Bellman equation is conceptually rather simple - % one can just use method of simple iteration: 
starting with some initialization $d_0(s)$, the new approximation for the distance function is iteratively updated by
by the right hand side of the Bellman equation (method is also called dynamic programming in RL context), i.e.:

% At the preliminary stage of the research we performed experiments with dynamical programming (DP) methods~\cite[Ch.~4]{sutton2018} for graphs small enough to fit into memory. Algorithm~\ref{alg:DP} is a DP version suitable for solving the Bellman equation~\eqref{bel_eq}. Starting with some initialization $d_0(s)$, the new approximation for the distance function is iteratively updated by

$$
d_k(s) = 1+\min_{t \in N(s)} d_{k-1}(t), \quad d_{k}(e) =0.
$$

\begin{algorithm}[H]
 \caption{Simple iteration (Dynamic programming)}
 \begin{algorithmic}[1]
 \STATE \textbf{Hyperparameters:} $\epsilon > 0$ (tolerance), $\alpha \in (0, 1]$ (learning rate), $N$ (maximum iterations)
 \STATE Initialize $d_0(s)$ for all states $s \in \mathcal S$ (e.g., by zeros)
 \FOR{$i = 0$ to $N$}
 \STATE $d_{i+1}(e) \leftarrow 0$
 \FOR{each state $s \ne e$}
 \STATE 
 $d_{i+1}(s) \leftarrow  \alpha(1 + \min\limits_a d_i(a\circ s)) + (1-\alpha)d_i(s)$
 \ENDFOR
 \STATE
 $c_i \leftarrow \mathrm{corr}(\boldsymbol d_i, \boldsymbol d)$
 \IF {$\Vert \boldsymbol d_i - \boldsymbol d\Vert  < \epsilon$}
 \STATE \textbf{break}
 \ENDIF
 \ENDFOR
 \STATE \textbf{Output:} Number of iterations $i$, vector of correlations $\boldsymbol c$ 
 \end{algorithmic}
 \label{alg:DP}
 \end{algorithm}

{\bf Examples of solving Bellman equation. Initialization dependence.}
Let us illustrate solving Bellman equation on Cayley graphs and demonstrate some unexpected dependence on initializations.
We observe that if initialization is positive convergence takes not more than diameter steps (it is no longer true for initialization with negative values). However convergence might not be monotonic - depending on initialization, it might happen that first
solution deviates from the final result, and only after several steps, begins to converge again. 

For small graphs such as \href{https://www.kaggle.com/code/fedmug/bellman-rc2}{Pocket Cube} or \href{https://www.kaggle.com/code/fedmug/rl-dynamic-programming-on-small-lrx}{LRX with $n\leqslant 10$} the true distances $d(s)$ can be easily calculated via BFS for each state. Therefore, we can use these true distances as stopping criteria: when $\|d_n(s) - d(s)\| < \varepsilon$, we are finished. In addition, we measure performance on each iteration by calculating the correlations between $d(s)$ and $d_k(s)$.

The main question of interest here was how the initialization influences the performance. That is important to understand in view of our proposal to use initialization coming from the diffusion distance.
We tried different initializations $d_0(s)$:

\begin{itemize}
 \item 
 zero initialization $d_0(s) = 0$ for all $s$
 \item 
 initialize each of $m$ layers by a specific value (e.g., first $k$ layers by true distances, last $m - k$ layers by zeros, $1 \leqslant k < m$)
 \item
 Manhattan distance to the target: 
 $d_0(s) = \sum \vert s_k - k \vert$
 \item 
 Hamming distance to the target:
 $d_0(s) = \vert \{k\colon s_k \ne k\}\vert$

\item 
 random integers form $0$ to $m$ or random values from $\mathcal N(0, \sigma^2)$

\item 
outputs of a model predicting diffusion distances (see next subsection or formula~\eqref{eq:dd})
 \end{itemize}

\begin{figure}[ht]
    \centering
    \includegraphics[width=1.0\linewidth]{figs/pearson.png}
    \caption{Pearson correlation between $d_i(s)$ and $d(s)$ on each iteration of the DP algorithm, $\epsilon = 10^{-3}$ (Pocket Cube)}\label{fig:pearson}
\end{figure}

These experiments revealed interesting phenomena presented in figure \ref{fig:pearson}. 
For some specific initialization the performance first degrades, and only after that achieves the convergence to the true distance which is guaranteed by the general theory. The number of steps for convergence does not exceed the diameter of the graph if initialization is positive, however, it can be larger in presence of negative values. Initialization coming from the machine learning model ("catboost") trained on diffusion distance (same as "warm-up step" in our proposal) improves the convergence, though not in a radical manner. 


{\bf Deep Q-learning algorithm.} Tabular methods like DP or Q-learning are not applicable in situations where the graph is too large. A family of deep Q-learning methods like~\cite{mnih2013playing} emerged to resolve this issue. Typically, these algorithms leverage a neural network to predict targets like right-hand sides of~\ref{eq:v-star-bellman} or~\ref{bel_eq}:

$$
    f_{\boldsymbol \theta}(s) = r + \max_a Q_{\boldsymbol \theta}(s, a).
$$

The algorithm~\ref{alg:deepQ} is a deep Q-learning method designed to solve the Bellman equation~\ref{bel_eq} (here we assume that $s$ is an embedding of a graph node suitable for input of the neural network $f_{\boldsymbol \theta}(s)$). Typically, the choice of the training subset of nodes is made by walking over the graph, combining both random steps and steps guided by the model. After the training is completed, the output of the model $f_{\boldsymbol \theta}(s)$ should hopefully serve as a good approximation to the true distance from $s$ to $e$.
%One assumes that a feature vector is attached to each graph node. 
%One considers a neural network that takes a feature vector as an input and should output the approximate solution of~\eqref{bel_eq}.
%The training process is organized as follows.
%For each epoch, one selects some subset of nodes, and from the current neural network, $d(s)$  computes the right-hand side of the equations \ref{bel_eq}
%$n(s) = 1+\min_{t \in N(s)} d(t), n(e) =0$, denoting %$n(s)$ as a result of that computation. 
%Afterwards, gradient descent is run to minimize the error between $d(g)$ and $t(g)$. If the error is zero, that means that $d(g)$ solves the Bellman equations. Thus, in graph theory setup, the distance function ("word metric" in the case of Cayley graphs) is produced.
%Typically, the choice of the training subset of nodes is made by walking over the graph, combining both random steps and steps guided by the model. We will discuss that in detail below.

\begin{algorithm}[t]
 \caption{Deep Q-learning}
 \label{alg:deepQ}
 \begin{algorithmic}[1]
 \STATE \textbf{Hyperparameters:} $N$ (number of epochs), $M$ (subset size)
 \STATE Initialize weights $\boldsymbol \theta$ of the neural network $f_{\boldsymbol \theta}(s)$
 \FOR{$i = 0$ to $N$}
 \STATE select a random subset of nodes $S_i$, $|S_i| = M$, $e\notin S_i$
 \FOR{each state $s \in S_i$}
 \STATE 
 $d_i(s) \leftarrow  1 + \min\limits_{t \in N(s)}f_{\boldsymbol \theta}(t)$
 \ENDFOR
 \STATE calculate MSE loss $\mathcal L(\boldsymbol \theta) = \frac 1M \sum\limits_{s\in S_i}\big(d_i(s) - f_{\boldsymbol \theta}(s)\big)^2$
 \STATE calculate $\nabla_{\boldsymbol \theta} \mathcal L(\boldsymbol \theta)$ and do the gradient step
 \ENDFOR
 \end{algorithmic}

 \end{algorithm}

{\bf Graph path finding.}
To find a path from a given node $s$ to the target  $e$, one needs to run a greedy search algorithm, which, starting from node $s$, inspects all its neighbors and evaluates distance function for them, then chooses the node with the minimum value of the distance function and makes a move to that node. The process is repeated until the number of the destination node found at the limit of the step is exceeded. If the distance function is the true distance on the graph (the shortest path length), such an algorithm finds the shortest path optimally (least possible operations). However, the predictions of the neural network $f_{\boldsymbol \theta}(s)$ from the deep Q-learning algorithm~\ref{alg:deepQ}, are typically not very precise, and the greedy algorithm will be stuck in local minima. To overcome that difficulty, one used more advanced graph search algorithms like beam search, A-star, etc. 

A guided beam search approach is described in algorithm~\ref{alg:beam_search}; on each iteration it maintains $W$ states with lowest estimations to the target state. If algorithm~\ref{alg:beam_search} can find the path to the target in $N$ iterations, it returns the length of this path.
%That is discussed below.


\subsection{Diffusion distance pathfinding (CayleyPy-1)}\label{sec-DD}
Let us briefly remind the alternative pathfinding approach proposed in our previous paper \cite{chervov2025machinelearningapproachbeats}. Its efficiency has been demonstrated on the Rubik's cube, where significantly outperforms all possible competitors, providing shorter solutions than any existing approach for $3\times 3\times3$, $4\times 4\times 4$, $5\times 5\times5$ cubes. In some sense, it is simpler and more computationally efficient.

\begin{algorithm}[t]
 \caption{Guided beam search pathfinding}
 \label{alg:beam_search}
 \begin{algorithmic}[1]
 \STATE \textbf{Hyperparameters:} $s_0$ (starting state), $M$ (step limit), $f(s)$ (distance estimator), $N(s)$ (neighbors of $s$), $W$ (beam width)
 \STATE $B \leftarrow \{s_0\}$
 \FOR{$i = 0$ to $M$}
    \STATE $B' \leftarrow \varnothing$
    \FOR{each state $s \in B$}
        \IF{$s = e$}
         \STATE \textbf{return} $i$
        \ENDIF
        \STATE $B' \leftarrow B' \cup N(s)$
    \ENDFOR
    \STATE $B' \leftarrow B' \setminus B$
    \STATE sort states $s \in B'$ by the distance estimations: 
    $$
    f(s_{(1)}) \le f(s_{(2)}) \le \ldots \le f(s_{(|B'|)})
    $$ 
    \vspace{-5mm}
    \STATE $B\leftarrow \{s_{(1)}, \ldots, s_{(\min\{W, |B'|\})}\}$
 \ENDFOR
 \end{algorithmic}
 \end{algorithm}
 

The idea is to work not with the true distance on the graph but with the diffusion distance, which is the length of the random path. The rigorous definition of the diffusion distance is given by the formula~\eqref{eq:dd} in the supplementary section.

% More precisely, for each state $s$ diffusion distance is defined as

% $$
%     \mathrm{DD}(s) = \frac{\sum\limits_{k=0}^{K_{\max}} k \mathbb P(s, k)}{\sum\limits_{k=0}^{K_{\max}}  \mathbb P(s, k)}
% $$

% \noindent
% where $\mathbb P(s, k)$ is the probability to reach $s$ from the initial state $e$ in $k$ steps. 
% %If $A$ is the adjacency matrix of the Cayley graph then $\mathbb P(s, k) = \frac{A^k_{0s}}{\sum\limits_j A^k_{0j}}$. 

The diffusion distance is easy to estimate by the number of steps of random walks. Generating random walks is much faster than computing neural network predictions. Thus, the training data generation is much more efficient than in Q-learning approaches, though the training data might be of poorer quality. 

%The assumption is that each graph node has a feature vector attached. For example, for permutation groups, it is a vector representing a permutation $(p_0,p_1,p_2,...p_{n-1}), p:i \to p_i $
%We also assume that there is a selected node (e.g., group identity) where we need to find a path from any other node. 

% The assumption is that for each node of a graph there is feature vector attached. For example for permutation groups it is a vector representing a permutation $(p_0,p_1,p_2,...p_{n-1}), p:i \to p_i $
% We also assume that there is a selected node (e.g. identity of the group $e$) where we need to find a path from any other node. 

The core of our approach consists of three steps:

\begin{enumerate}
    \item {\bf Generation of the training set via random walks.} 
    Generate $N$ random walk trajectories starting from the selected node $e$. Each random walk trajectory consists of up to $K_{\text{max}}$ steps, where $N$ and $K_{\text{max}}$ are integer parameters of the method. If a random walk visits node $s$ for the first time at step $k$, we store the pair $(s,k)$ in our training set. 
    %For some nodes encountered during the random walks, we store a set of pairs $(v, k)$, where $v$ represents the vector corresponding to the node and $k$ is the number of steps required to reach it via the random walk. 
    
    \item {\bf Training a model.}  
    The generated set of pairs $(s,k)$ serves as the training data for a supervised machine learning model $f(s)$, typically a neural network.
    %The generated set of pairs $(v, k)$ serves as the training set for a supervised model of machine learning (usually a neural network). %Specifically, $v$ serves as the 'feature vector' (the input for the neural network), and $k$ represents the 'target' (the output that the network needs to predict). Thus, 
    For a given node $s$ this model predicts the estimated diffusion distance from $s$ to the target $e$.

    \item {\bf Guided beam search.} Here we apply the algorithm~\ref{alg:beam_search} to find a path from a given state $s$ to the target node $e$. The model $f(s)$ trained on the previous step provides heuristics on where to make the next steps, while the beam search technique compensates for any possible incorrectness in the predictions of the neural network. Having fixed a positive integer $W$ ("beam width" or "beam size"), on each step we select $W$ nodes that are closest to the destination according to the neural network. Then we check $f(s)$ for all neighbors of these best $W$ states, drop duplicates, and repeat the process until the destination node is found or the step limit is exceeded.
    %This step finds a path from a given node to the destination node. The neural network provides heuristics on where to make the next steps, while the graph pathfinding technique compensates for any possible incorrectness in the neural network predictions. The beam search pathfinding method is quite simple, but has proven to be the most effective for us and works as follows. Fix a positive integer $W$ (called "beam width" or "beam size"). 
    %Starting from a given node, we take all its neighboring nodes and compute the neural network predictions for all of them. We then select the $W$ nodes that are closest to the destination according to the neural network (i.e., the predictions have smaller values). For these selected $W$ nodes, we take their neighbors, drop duplicates, and again compute the neural network predictions, choosing the top $W$ nodes with the best (i.e., minimal) predictions. This process is repeated until the destination node is found (or limit of steps is exceeded).
  
\end{enumerate}

The last step is common for any method, whether it is based on Q-learning or diffusion distance. One uses a neural network predictions as a heuristic function and graph search algorithms to find a path. We are strongly in favor of beam search since it is quite simple, but has proven to be the most effective for us, while the DeepCube team~\cite{mcaleer2019solving, agostinelli2019solving, khandelwal2024towards,agostinelli2024q} uses the A-star algorithm and develops its modifications.

A simple but useful modification of the diffusion distance method leverages \textbf{non-backtracking} which is applied during steps~(1) and~(3). On both occasions, we store the history states visited in the $m$ previous steps to avoid returning to them. The history of non-backtracking random walks contains $O(mN)$ states, and all trajectories are not allowed to visit any of these states and their neighbors in the next step. For non-backtracking beam search we do not include the states from the history on line 9 of the Algorithm~\ref{alg:beam_search}.

\subsection{Beam search effectivity. Toy model of research community }\label{sec-Beam}

To some extent, the success of the proposed above is due to beam search and our efficient implementation supporting large beams. We started with a modified greedy search, then moved to Metropolis and A-star, before finally finding relief in beam search. The larger the beam size, the more effective beam search becomes (probability to find path increases and paths become shorter). To maximize beam sizes, we developed an original and efficient PyTorch implementation capable of supporting beam sizes with millions of nodes and beyond (\href{https://www.kaggle.com/code/alexandervc/beamsearch-basic}{basic}, \href{https://github.com/iKolt/cayleypy/blob/main/cayleypy/cayley_graph.py#L447}{advanced} versions). On one hand, beam search is arguably one of the simplest methods. On the other hand, there are several key aspects worth highlighting.

The primary reason we need beam search is to avoid the problem of local minima. Greedy search can get stuck in them, whereas beam search can bypass them entirely if the beam size is larger than the depth of the local minima. It is similar to how an elephant wouldn't even notice small dips in the terrain, while for an ants, they could be an insurmountable obstacles.

Another important aspect of beam search is that increasing the beam size increases both computation and memory consumption linearly. Typically, memory -- not computation time -- is the main bottleneck. However, maximizing the beam size is crucial for performance.
%important aspect of beam search is that increasing the beam size not only linearly increases computation but also linearly increases memory consumption. On one hand, typically memory - not computational time - is the main bottleneck of the method. On the other hand, having the largest possible beam size is crucial for its performance.
Moreover, there is no simple way to reproduce the results of beam search with a beam size of $B+1$ using any algorithm that only consumes memory equivalent to beam size $B$.

To explain this point, we will use the analogy with the work of the research community.
Imagine two researchers working together on a creative problem (a beam size of 2). The key point is the exchange of ideas: if one researcher makes a breakthrough, they share it with the other, allowing both to continue their search from the breakthrough point. Without this exchange, the second researcher might get stuck in a dead end. This sharing of ideas helps avoid wasting time in unproductive directions.
The mechanism described above explains why increasing the beam width is so effective: it acts as an exchange of ideas, focusing on breakthroughs rather than following unproductive paths.
Essentially, this is what beam search does: it selects the top nodes across the entire beam neighborhood rather than picking the best few for each individual node. This means that if the neighbors of one node are significantly closer to the destination (analogous to a breakthrough), those nodes are retained in the beam, while the neighbors of less successful nodes are discarded. In other words, the descendants of some nodes do not survive into the next generation. A kind of natural selection process.

%From the above, it is clear that the size of a research community matters. Reducing the community size by half will not simply cut the research output in half -- it will decrease it exponentially. That aspect is different from the more routine work - cutting it twice output will be cut twice. 

These observations show that community size has a disproportionately large impact on research productivity. Halving the number of researchers does not simply halve the output — it causes a super-linear (often exponential) decline. In contrast, for routine tasks, a 50\% reduction in resources usually yields about a 50\% reduction in output. The same principle applies to beam search: the beam size is crucial, and there is no simple way to achieve the performance of a beam size $B$ using a smaller beam.

This logic also explains why evolutionary or nature-inspired global optimization methods can sometimes be successful. They rely on the same fundamental mechanism as beam search but offer better control over the diversity within the beam.
In some cases, it is important not to immediately jump to breakthroughs made by others but to continue exploring independent directions, as this might lead to even greater discoveries. For example, if medical doctors would stop doing their job to start follow breakthroughs in mathematics, it would lead to disastrous results — diversity is crucial in the long run.
The balance between "everyone follows the latest breakthrough" (as in beam search) and maintaining diversity is delicate. There is no universal solution, as different tasks require different trade-offs. Exploring various methods on Cayley graph pathfinding tasks is an interesting direction.

\subsection{Combining diffusion distance with deep Q-learning} \label{sec-RLDD}
%~\\
%{\bf Rationale. Why we need that?} 
The diffusion distance approach described above is quite simple and efficient; however, there are practical and theoretical disadvantages. The practical side consists of somewhat surprising phenomena that enlarging the training set (generating more random walks) does not always improve the quality of the model; such stagnation was observed for Rubik's cube cases in our previous paper, and it is even more prominent for LRX graphs. It is actually related to the quite clear theoretical fact: diffusion distance is not monotonic function of the true distance.
%one node can be further than another for the true distance, while oppositely ranged for the diffusion distance. 
This problem is not only for random walks estimation, but also for the diffusion distance itself (e.g., see computed theoretically \href{https://www.kaggle.com/code/fedmug/lrx-diffusion-distance-analysis}{analysis}).
%(i.e. limit of infinite data - computed theoretically (\href{https://www.kaggle.com/code/fedmug/lrx-diffusion-distance-analysis}{analysis}) ). 
%Such phenomena will be discussed below in more details (see also our n with analysis). 
So, diffusion distance is easy to estimate, but it is not precise. On the other hand, DQN and Bellman's equation allow us to approximate the true distance, but for the price of more heavy computations. 
The idea of the proposal below merges lightweight diffusion distance with a heavier and more precise Bellman equation approach, and at avoiding the sparse reward problem.


Our proposal looks as follows.

{\bf Part 1. Warm-up diffusion distance training.} Apply diffusion distance approach for pre-training the neural network, i.e. steps (1) and (2) above: generate training data via random walks, and train a neural network $f_{\boldsymbol \theta}(s)$ to predict the number of steps of the random walk.

{\bf Part 2. Modified DQN training.} We use the deep Q-learning algorithm~\ref{alg:deepQ} with two enhancements:

\begin{itemize}
    \item 
    the random subset of nodes consist of all states visited by $N$ random walks of size $K_{\max}$;
    \item 
    the result of the Bellman update $1 + \min\limits_{t \in N(s)}f_{\boldsymbol \theta}(t)$ is clipped to the interval $[0, k]$ where $k$ is the time step when the state $s$ was visited for the first time.
\end{itemize}

The pseudo code for this modified version of DQN is presented in Algorithm~\ref{alg:deepQ_modified}.

%The second phase relies on the Bellman equation strategy with mild modifications and employs a somewhat different space exploration approach compared to conventional DQN. This strategy is inspired by the diffusion distance method:

%Step 1. Generate $N$ non-backtracking random walk trajectories, all starting from the target node $e$. Each random walk consists of up to $K_{\text{max}}$ steps, where $N$ and $K_{\text{max}}$ are integer parameters of the method. 

%Step 2. For the nodes visited during the walks, compute Bellman equation-based predictions:
%$t(g) = 1+\min_{n:~neigbours~of~g} d(n), t(e) =0$ %using the standard formulation. 

%Step 3. Clip the predictions based on the number of steps taken in the random walk, since the true distance cannot exceed the number of random walk steps. Also, clip any negative predictions to zero.

%Step 4. Run gradient descent to train the neural network by minimizing the loss function between the clipped $t(g)$ values (as targets) and the neural network's predictions. 

{\bf Part 3. Guided beam search.} Here we apply step (3) from the diffusion distance approach described above (Algorithm~\ref{alg:beam_search}).

%{\bf Part 3. Graph search algorithm - beam search with a neural network as a heuristic. } After the neural network is trained, we utilize the same approach to find a path using a neural network-guided beam search as in our previous paper.

\begin{algorithm}[t]
 \caption{Modified deep Q-learning}
 \label{alg:deepQ_modified}
 \begin{algorithmic}[1]
 \STATE \textbf{Hyperparameters:} $M$ (number of epochs), $N$ (number of random walks), $K_{\max}$ (length of random walks), $f_{\boldsymbol \theta}(s)$ (distance estimator)
 \FOR{$i = 0$ to $M$}
 \STATE Generate $N$ random walks of size $K_{\max}$ from the target node $e$
 \STATE Collect dataset $\mathcal D = \{(s, k)\}$ of states $s$ visited by random walks at time $k$
 \FOR{each $(s, k) \in \mathcal D$}
  \STATE 
  $d(s) \leftarrow  1 + \min\limits_{t \in N(s)}f_{\boldsymbol \theta}(t)$ (update based on Bellman equation~\eqref{bel_eq})
  \STATE $d(s) \leftarrow \min\{k, \max\{0, d(s)\}\}$ (clipping)
 \ENDFOR
 \STATE Calculate MSE loss $\mathcal L(\boldsymbol \theta) = \frac 1{|\mathcal D|} \sum\limits_{\mathcal D}\big(d(s) - f_{\boldsymbol \theta}(s)\big)^2$
 \STATE Calculate $\nabla_{\boldsymbol \theta} \mathcal L(\boldsymbol \theta)$ and do the gradient step
 \ENDFOR
 \end{algorithmic}

 \end{algorithm}

There are several advantages of the proposed method. First, it avoids the sparse reward problem: since we always start from the selected state, we can reliably clip overestimated predictions. Second, the warm-up phase helps avoid the long initialization typical for standard DQN, where targets from the Bellman equation are noisy because the network starts with random weights. In our approach, the targets are meaningful from the beginning: in the first phase, they come from random walk steps; in the second phase, the network itself produces useful targets, as it has already been pre-trained. Finally, we use non-backtracking random walks for a clear reason: the number of steps in such walks better reflects the true distance. Conceptually, this relates to the well-known phenomenon "Non-backtracking random walks mix faster"~\cite{alon2007non}.

%\section{Benchmarks}
\section{Analysis and results}

%Goal - $S_{110}$. 
%DeepMind's success with Go game with $10^{170}$ states - creates a challenge to achieve similar results in the other fields in particular for Cayley graphs, at the moment, we are somewhat far from that, but there are plenty ideas not yet incorporated in our CayleyPy project, and optimistically such a goal is achievable in near horizons.
DeepMind's success in mastering the game of Go, which has approximately $10^{170}$ possible states, poses an inspiring challenge: achieving similar breakthroughs in other domains, such as Cayley graphs. At present, we are still a distance from this goal. However, many promising ideas have yet to be implemented in our CayleyPy project, and with continued development, reaching such a milestone may be possible in the near future.

In this study, we demonstrate successful pathfinding in very large LRX Cayley graphs of $S_n$. We use its conjecturally longest element~\eqref{eq:longest} of expected length $n(n-1)/2$ as a benchmark.  Our zero prior knowledge approach can find the paths from $\ell_n$ to $e$ for $n \leqslant 44$. This is a significant improvement over the performance of the GAP computer algebra system, which can only handle up to around $n=20$, and with much longer path lengths and runtimes.  Introducing some domain knowledge (see Section 3.2 below) drastically improves the performance to $n=100$ or even more.
%However, so far, $n=40$ remains out of reach, despite extensive experimentation with different neural network architectures. Surprisingly, nearly all approaches we tested stuck slightly above $n=30$ showing only minor variation. 

%But somewhat far from the desirable goal

\subsection{"Overcome the GAP" with zero prior knowledge AI }
Here we demonstrate that the diffusion distance method significantly outperform classical methods of computer algebra system GAP. 

\begin{table}[ht]
    \centering
    \begin{tabular}{|c|c|c|c|c|c|}
        \hline
        $n$  &  GAP length   & $\frac{n(n-1)}2$            &   DD length    \\
        \hline
        9    &  41     & 36               &  36   \\
        \hline
        10   &    51   & 45               &  45       \\
        \hline
        11   &     65      & 55             &  55     \\
        \hline
        12   &    78       &  66          & 66     \\
        \hline        
        13   &    99       &  78           &   78       \\
        \hline        
        14   &   111       &  91           &   91     \\
        \hline        
        15     &   268      &  105         &  105    \\
        \hline        
        16     &    2454    &  120         &  120     \\
        \hline        
        17   &     380     &  136         & 136    \\
        \hline        
        18   &     20441   &   153        & 153    \\
        \hline        
        19   &      3187      & 171        & 171    \\
        \hline      
        20   &     217944    & 190        & 190      \\
        \hline 
        21   &   -           & 210     & 210        \\
        \hline           
    \end{tabular}
    \captionsetup{skip=10pt} 
    \caption{Comparison of GAP and diffusion distance (DD) method for the conjecturally longest element of LRX Cayley graph with expected length $n(n-1)/2$.}
    \label{tab:example}
\end{table}

For $n=20$, the  GAP timing is 41min 18s,  while our methods can find results much faster.
For example, the diffusion distance method with simple neural network requires only 3m 48s of calculations on GPU P100 (\href{https://www.kaggle.com/code/alexandervc/lrx-cayleypy-rl-mdqn?scriptVersionId=224270083}{notebook version 423}). GAP benchmarks have been performed \href{https://www.kaggle.com/code/avm888/group-elements-decomposition-gap-limits}{here}. 


\begin{table}[ht]
    \centering
    \begin{tabular}{|c|c|c||c|c|c|}
\hline
$n$  & $\frac{n(n-1)}{2}$ & DD length &
$n$  & $\frac{n(n-1)}{2}$ & DD length \\
\hline\hline
22 & 231 & 231  & 31 & 465 & 478 \\
\hline
23 & 253 & 253  & 32 & 496 & 508 \\
\hline
24 & 276 & 276  & 33 & 528 & 642 \\
\hline
25 & 300 & 300  & 34 & 561 & 667 \\
\hline
26 & 325 & 325  & 36 & 630 & 834 \\
\hline
27 & 351 & 351  & 38 & 703 & 759 \\
\hline
28 & 378 & 378  & 40 & 780 & 926 \\
\hline
29 & 406 & 408  & 42 & 861 & 909 \\
\hline
30 & 435 & 437  & 44 & 946 & 1060 \\
\hline
\end{tabular}
    \captionsetup{skip=10pt} 
    \caption{Lengths of solutions produced by diffusion distance (DD) method for the conjecturally longest element of LRX Cayley graph with expected length $n(n-1)/2$.  }
    \label{tab:example_nn}
\end{table}

The \texttt{DD length} columns in Tables~\ref{tab:example} and~\ref{tab:example_nn} are aggregated results from experiments with different guiding models based on gradient boostings and neural networks.
All of them solve cases with $n\le 28$ quite stably; moreover, after some hyperparameter tuning, they find the shortest paths. Going beyond the limit $n=28$ requires some effort. For $n\le 44$ a path can be occasionally found by multiple relaunches of the pipeline, however, its length is no longer optimal.
%, and in several steps in $n$ they loose ability to find a path. 
%It is surprising that all types of models have the problem at the same range of $n$ and that length is not becoming gradually longer and longer, but  it  stays almost perfect, and then  quite sharply pipeline loses ability to find a path at all. 
%For some cases  we use several reruns - 
%Occasionally, if the path was not found via a single run, it may sometimes be found by multiple relaunches of the pipeline. 
%For example, case $n=32$ has been solved with our modified DQN strategy with $30$ warm-up epochs and $1000$ DQN epochs (\href{https://www.kaggle.com/code/rustemturtayev/lrx-cayleypy-rl-mdqn-d44272?scriptVersionId=223961901}{notebook}), time: 2h 53m 53s GPU P100. 

\subsection{X-trick.} Surprisingly, adding one line of code allows us to drastically increase the size of solvable graph, from $S_{40}$ to $S_{100}$. Namely, add the following condition in line~9 of the beam search algorithm~\ref{alg:beam_search}: if $s_0 < s_1$ then do not include $X\circ s$ in $N(s)$.
%\subsection{One line of code to get orders of magnitude improvement from $S_{30}$ to $S_{100}$. }
%If the following condition is applied in the beam search:
%\texttt{if action == 'X' and current\_state[0] < current\_state[1]: continue }
%then a dramatic improvement in performance is observed. 
The logic is straightforward: if the first two elements of the permutation are already sorted, it makes little sense to swap them by applying the generator $X = (0 \ 1)$. 
%- if the contents of positions 0 and 1 are already sorted - it is logical not to swap them i.e. not to apply $X$ generator. 

This X-trick turns out to be a game changer. By adding it to the pipelines that succeeded for $n$ about $40$, we were able to find paths for $n=100+$ (and potentially even larger values, as we were limited by the 16GB RAM limit). This represents a shift from working with Cayley graphs of size $10^{30+}$ to $10^{150+}$.  This trick worked universally across model types, including perceptrons, gradient boosting machines, and convolutional neural networks. The key challenge now is how to eliminate such prior knowledge and generalize the method to other types of generators. However, X-trick may prevent the discovery of optimal paths, especially for small values of $n \le 15 $.


Examples of lengths found by the diffusion distance method with X-trick are presented
in Table~\ref{tab:lens_upto100} below:

\begin{table}[ht]
    \centering
    \begin{tabular}{|c|c|c|c|c|c|}
        \hline
        $n$   & 100 & 90   & 80  & 70 & 60  \\
        \hline
        Solution length        & 16632 & 33459  & 7510  & 8115 & 3650 \\
        \hline
        Diameter (proposed) &  4950   & 4005  & 3160  & 2415 & 1770 \\
        \hline
        Time         &  8h 43m & 3h 57m   & 1h 5m  & 52m & 3h  \\
        \hline
        Notebook version   &   255   &  242     & 241    & 233  & - \\
        \hline
    \end{tabular}
    \captionsetup{skip=10pt} 
    \caption{Lengths of solutions for different $n$ obtained by ML model followed by beam search with X-trick. (Not completely zero prior knowledge approach). }
    \label{tab:lens_upto100}
\end{table}

Computation for $n=100$ took 8h 43m 53s GPU P100. Beam width $ {2}^{20}$, 30 epochs of training ( \href{https://www.kaggle.com/code/alexandervc/lrx-cayleypy-rl-mdqn?scriptVersionId=223954525}{Notebook version 255}). For other cases we used beam size $\le {2}^{18}$, that is why for $n=90$ length is much larger. The data for $n=60$ is obtained by the gradient boosting model, with slightly better lengths, but more timing, in the other cases single layer MLP model used. In all cases we decompose the conjecturally longest element of expected length $n(n-1)/2$. 

\subsection{Diffusion distance + modified DQN}
We observe that replacing the training stage with the DQN algorithm~\ref{alg:deepQ_modified} gives a consistent improvement over the diffusion distance approach; however, the effect is not very large. Below we provide two illustrative experiments.

\begin{table}[ht]
    \centering
    \begin{tabular}{|c|c|c|c|c|}
        \hline
        $n$  & epochs: $5+0$  & epochs: $50+0$   & epochs: $5+50$   & epochs: $50+50$ \\
        \hline\hline
        6 &  0.955 & 0.99  & 0.99  & 0.99 \\
        \hline
        8 &  0.954 & 0.969 & 0.98   & 0.982 \\
        \hline
        10 & 0.883  & 0.923 & 0.939   &  0.938 \\
        \hline
    \end{tabular}
    \captionsetup{skip=10pt} 
    \caption{Spearman correlation between the true distances and the predictions of the neural network trained for $A$ epochs during the warm-up diffusion distance phase and $B$ epochs during the DQN phase (epochs: $A+ B$). Additional training with DQN leads to improved results.}
    \label{tab:DQN_small}
\end{table}

The first experiment (see Table~\ref{tab:DQN_small}) focuses on relatively small groups, where the true graph distance can be computed by brute force and compared with the neural network's predictions. For small group sizes, the method learns the graph distances almost perfectly. However, for larger groups, the correlation between predictions and true distances becomes less accurate. Still, additional training using the modified DQN always improves upon the initial warm-up phase based on the diffusion distance procedure.  

\begin{table}[ht]
    \centering
    \begin{tabular}{|c|c|c|c|c|c|}
        \hline
        Warm-up epochs  & DQN epochs  & success rate & shortest path & median path \\
        \hline\hline
        30 & 200 & 80\%  &  382 & 402  \\
        \hline
        30 &  0 & 70\%  & 464 &  650  \\
        \hline
        100 &  0 & 70\%  & 436 & 516  \\
        \hline
        0 &  230 & 20\%  & 434 & 471  \\
        % 2 &  396 & 658   & 516  &  None \\% AC
        % 3 &  382 & None  & None  & None  \\% AC
        % 4 &  408 & 650   & 436  &  None \\ % AC
        % 5 &  428 & 706   & 1046  & None  \\ % Fironov
        % 6 &  396 & None  & None  & 508  \\ % Fironov
        % 7 &  426 & 890   & 614  &  None \\ % Fironov
        % 8 &  None& 598   & None &  None \\ %Rustem
        % 9 &  392 & 464   & 524  &  None  \\ % Shishina
        % 10 & None& 626   & 452  &  434 \\ %Liuda 
        % Median & 402& 650   & 516 & 455.5  \\
        \hline
    \end{tabular}
    \captionsetup{skip=10pt} 
    \caption{Path lengths for decomposing the longest element in the LRX Cayley graph of $S_{28}$.
    Results in each row are based on 10 independent runs. Additional training with DQN improves both the average path length and the success rate of finding a path. Pure DQN training alone performs poorly. }
    \label{tab:DQN_28}
\end{table}

The second experiment concerns the larger group $S_{28}$ with LRX generators. We train the model and use beam search to find a path, where the path length serves as the performance metric. The results provided in Table~\ref{tab:beam_size} confirm that additional training with DQN improves performance. The combination of $30$ warm-up and $200$ DQN epochs produces one of the best outcomes observed in this setting. The median (shortest) path length $402$ ($382$) is very close to the expected ideal length of $378$. For each run, models were retrained from scratch; thus, due to the randomness of the walks, the resulting models differ significantly. The analysis can be found in the Kaggle \href{https://www.kaggle.com/code/alexandervc/lrx-cayleypy-rl-mdqn}{notebook}. 

\subsection{Non-backtracking}
We exploit non-backtracking idea in two ways - first we can use non-backtracking random walks, and the second we can use non-backtracking in beam search. The second way shows quite a good effect. 
% The influence of non-backtracking random walks is in a way ambiguous. On the one hand, for some values of $n$, introducing nonbacktracking allows finding a path from the longest element $\ell_n$ to $e$, even if pure random walks completely fail. On the other hand, the performance seems to be almost independent of the history size $m\geqslant 2$ (see Table~\ref{tab:nbt}).
%Non-backtracking random walks have quite a strong effect. 
%We actually use improved version of the non-backtracking based on the following ideas: 

% \begin{enumerate}
%     \item
%  we forbid not only 1 step in history but typically 32 steps (it is a parameter);
%     \item
%     we consider a bunch of trajectories and ban states for all of them at once;
%  \item
%  we ban not only visited states but their neighbors.
% \end{enumerate}

%The price for that is slowing down the generation, but that is not so important since gradient descent takes a much longer time, and overall, there is no significant slowdown. 

\begin{table}[t]
    \centering
    \begin{tabular}{|c|c|c|c|}
        \hline
        history size, $m$ & success rate & median length & min length \\
        \hline\hline
        % $0$    &       $0.00\%$    & None     & None   \\
        % \hline
        2   & $40.41\%$ & $476$ & $380$ \\
\hline
4   & $40.10\%$ & $474$ & $382$ \\
\hline
8   & $40.32\%$ & $471$ & $386$ \\
\hline
16  & $37.40\%$ & $470$ & $382$ \\
\hline
32  & $37.47\%$ & $470$ & $378$ \\
\hline
40  & $36.16\%$ & $480$ & $390$ \\
\hline
48  & $37.39\%$ & $465$ & $380$ \\
\hline
128 & $36.45\%$ & $476$ & $382$ \\
\hline
254 & $37.85\%$ & $468$ & $382$ \\
\hline
400 & $37.40\%$ & $474$ & $384$ \\
\hline
512 & $36.58\%$ & $476$ & $384$ \\
\hline    
    \end{tabular}
    \captionsetup{skip=10pt} 
    \caption{Performance of the DD method with non-backtracking random walks ($n=28$). }
    \label{tab:nbt}
\end{table}

\begin{table}[ht]
    \centering
    \begin{tabular}{|c|c|c|c|c|c|}
        \hline
        Beam   &  Solution Length                  \\
       Width   &  Non-backtracking RW    \\
        \hline
        $2^{10}$    &       None           \\
        \hline
        $2^{12}$   &       252             \\
        \hline      
        $2^{14}$    &       202             \\
        \hline    
        $2^{16}$    &       194            \\
        \hline 
        $2^{18}$    &       194             \\
        \hline  
        $2^{20}$     &       190 (Ideal)    \\
        \hline        
    \end{tabular}
    \captionsetup{skip=10pt} 
    \caption{Influence of beam width and non-backtracking random walks for solving the longest element of the LRX Cayley graph $S_{20}$.     }
    \label{tab:beam_size}
\end{table}

%We generate 10000 trajectories at once each of the length 190, by the generalized nonbacktracking procedure above and train the neural network with batch size 1024 just for single epoch.
%The total time for each computation (both training and beam search stages) is 1--5 minutes; experiments can be found in versions 427--430 of the \href{https://www.kaggle.com/code/alexandervc/lrx-cayleypy-rl-mdqn?scriptVersionId=224294767}{notebook}.

The theoretical reason why non-backtracking is helpful is clear for the random walks. Imagine our graph is a tree, then  for a non-backtracking  random walk, the number of steps coincides with the true distance (length of the shortest path). Thus diffusion distance coincides with  the standard
distance. For complicated graphs, this is not the case, but still, it forces the diffusion distance to be closer to true distance and to have less variability. Here by variability we mean the situation when the same node can be achieved in a different number of steps. This is equivalent to having noise in the training data which makes models less precise. Similar for beam search non-backtracking mitigates stagnation
in local minimums. 

\begin{table}[ht!]
\centering
\begin{tabular}{|c|c|c|c|}
\hline
history size $m$ & success rate & median length & min length \\
\hline
2     & $17.65\%$ & $430$ & $384$ \\
\hline
4     & $17.08\%$ & $432$ & $390$ \\
\hline
8     & $21.37\%$ & $440$ & $382$ \\
\hline
16    & $30.07\%$ & $446$ & $388$ \\
\hline
32    & $38.55\%$ & $473$ & $384$ \\
\hline
64    & $43.29\%$ & $495$ & $380$ \\
\hline
128   & $56.95\%$ & $517$ & $384$ \\
\hline
256   & $60.61\%$ & $588$ & $384$ \\
\hline
512   & $67.33\%$ & $582$ & $384$ \\
\hline
1024  & $69.71\%$ & $663$ & $382$ \\
\hline
\end{tabular}
\caption{Performance of the DD method with non-backtracking beam search ($n=28$). The path is found more often as the history size grows. }
\label{tab:beam_ban}
\end{table}

Injecting non-backtracking into the beam search stage (Algorithm~\ref{alg:beam_search}) allows to increase the success rate: the greater is the number of banned states not included in the beam, the higher is the probability of finding a solution (see Table~\ref{tab:beam_ban}). On the other hand, the median solution length tends to increase when the history size is large. But that cannot be interpreted as performance degradation, rather it means that pipeline can solve more difficult cases, and hence getting large lengths.  


\subsection{Architectures and parameter dependence. }
%The tests of different architectures and parameters is discussed below.
%The main conclusion is that 
We observe that rather simple models like MLP and gradient boostings currently perform better than more advanced ones: transformers, CNNs, etc. An exception is one of our transformer models that can find paths up to $n=100$, however, it is not a zero prior knowledge model since it uses features specific to LRX generators. 

The dependence on various parameters was explored and discussed in our previous paper ~\cite{chervov2025machinelearningapproachbeats} for the case of the Rubik's cube group. The beam size proved to be one of the most important parameters: the length of solution is almost linearly improving on logarithm of beam size. Another observation is that it is more important for a neural network to be deep than wide, however, the effect of depth is better seen for large groups. We also noted the performance stagnation as size of the training set increased.  

For the LRX Cayley graph these phenomenons are less strongly expressed, or even negligible. 
Most experiments with deep neural networks have not lead to significant improvements over simple one layer perceptron with 128 neurons. In particular, one can train models in just a single epoch within a minute, e.g., for $S_{20}$. The beam size effect is still present but not as pronounced (see table~\ref{tab:beam_size}).
%Beam size effect is still clear, but less expressed (see table~\ref{tab:beam_size}). 


{\bf CNN.}
Convolutional neural networks have demonstrated a slower training process. Their results are similar to those of perceptrons, e.g., length of the solution for $n=16$ is 122 (ideal is 120), for $n=20$ is 216 (ideal is 190). The basic version of the architecture uses two 1-d CNN layers with batch normalization and dropout to prevent overfitting. One of the successful variations is using temporal gated CNN. The first part of the architecture is the residual block - several convolutional blocks with gating based on sigmoid activation function. The second part of the architecture is stacking several layers of residual blocks with gradually increased dilation to increase the size of the receptive field.
Notebooks: \href{https://www.kaggle.com/code/artgor/encoding-as-permutation-matrix-1d-cnn}{1}, \href{https://www.kaggle.com/code/artgor/encoding-as-permutation-mat-simple-2d-cnn}{2}, \href{https://colab.research.google.com/drive/1oWuQg33bveHYqVET4IES_EWIH1nsa4yi}{3},\href{https://colab.research.google.com/drive/1oWuQg33bveHYqVET4IES_EWIH1nsa4yi}{4}.

{\bf Transformer.} We also performed experiments with a transformer model. However, it was able to solve only not-so-large groups like $n\le 15$. 
That model incorporated standard transformer components, including multi-head self-attention, positional encoding, feed-forward layers, and dropout to mitigate overfitting. However, the model struggled to converge to a satisfactory solution, even for relatively short sequences. This issue became increasingly severe as the sequence length grew. For instance, with $n=9$ (corresponding to path length = 107 and ideal = 36) or $n=13$ (path length = 332, ideal = 78), increasing the projection dimension to 512 or 1024, adding more layers, or training for more epochs failed to address the underlying challenges. Notebooks: \href{https://www.kaggle.com/code/nursmen/lrx-transformer-training}{1}, \href{https://www.kaggle.com/code/nursmen/lrx-transformer-use}{2}.


{\bf Transformer based on LRX specific features.}
The features inspired by Hamming distance were used:  difference between an element and its position;  difference with the left neighbor;  difference with the right neighbor.  A standard multihead attention with one encoder layer and MLP was used to estimate the number of steps required to reach the target state. The transformer demonstrated the ability to handle sequences up to $n=15$ without using beam search. For larger $n$, beam search is employed, and the maximum $n$ for which results were obtained is 100. The model and features only indirectly depend on $n$, therefore, when trained on micro batches constructed as random walks from the identity permutation for $n(n-1)/2$ moves, it can be used for different values of $n$, both larger and smaller.  (\href{https://www.kaggle.com/code/sergeifironov/permutation-transformer-solver-100}{Notebook}). 

{\bf Gradient boosting.} We also applied CatBoost and XGboost in our diffusion distance pipeline. Surprisingly, they perform not worse than neural networks (it was not the cases for Rubik's cube where boostings were not able to solve $4\times 4 \times 4$, $5\times5 \times5$ cubes), and even produce a bit shorter solutions. 
(\href{https://www.kaggle.com/code/antoninadolgorukova/lrx-gbm-powered-pathfinding}{Code}).

More simulation results for various parameters can be found in \href{https://www.kaggle.com/competitions/lrx-oeis-a-186783-brainstorm-math-conjecture/discussion/612338}{spreadsheet here}, or  on \href{https://docs.google.com/spreadsheets/d/1KFYTqPDSdH8_vdvWz3YLTLKg5IlW1Kkx_9kJClSKGOs/edit?usp=sharing}{google sheets}.

\section{Mathematical contributions}

In this section, we prove lower and upper bounds for the diameter and for particular elements, describe the decomposition of the conjecturally longest element of the desired length $n(n-1)/2$, describe algorithms to decompose in LRX generators, and formulate several conjectures on the growth, random walks and spectrum of the LRX Cayley graph. 

\subsection{Long and the longest element}

\begin{figure}[ht!]
   \centering
   \includegraphics[width=1.0\linewidth]{figs/LongEl.jpg}
   \caption{Permutations arising from dihedral symmetries are provably long elements. The longest is conjecturally one of them.}\label{fig:LongEl}
\end{figure}

We created an efficient implementation of the brute force breadth first search algorithm (see this \href{https://www.kaggle.com/code/ivankolt/lrx-4bit-uint64}{notebook}) using bit operations over \texttt{int64} and were able to traverse all the permutations for $n \le 15$, outperforming OEIS data by two steps (for $n=15$ computation took several days on Hadoop cluster). From that computation we observed that the longest element is unique and has a specific form~\eqref{eq:longest}. It can be described as a symmetry of a regular polygon (Fig.~\ref{fig:LongEl}), which is orthogonal to 0--1 edge. We conjecture that this element is the only longest element with length $n(n-1)/2$. 
Surprisingly, the full inversion $(n-1, n-2, \ldots, 2, 1, 0)$ is not the longest. As we will prove below, actually all the permutations which correspond to symmetries of a regular polygon are also long elements.  

\clearpage

\begin{proposition}
    Let 
    $$
    n > 2,\quad m = \Big\lfloor \frac n2 \Big\rfloor,\quad 
    \delta_n = [n = 2m] = \begin{cases}
        1, n \text{ \rm{is} even},\\
        0, n \text{ \rm{is} odd}.
    \end{cases}
    $$
    The element $\ell_n = (1\ 0 \ n-1 \ n - 2 \ldots 3 \ 2)$
    has the following explicit decomposition into product of LRX generators:
\begin{equation}
\label{eq:long}
\ell_n = \prod_{i=1}^{m-\delta_n} \big( X L^{(-1)^{i-1}} \big)^i \circ \prod_{i=m -1}^1 \big(L^{(-1)^{i-\delta_n}}X \big)^i \circ R^m.
%\prod_{i=1}^{[ (n-1)/2]} \big( X L^{(-1)^{i-1}} \big)^i 
%\prod_{i=([ n/2]-1)}^{1} (L^{(-1)^{i+[(n-1)/2] - [n/2] }} X )^i R^{[n/2]}
\end{equation}    
\end{proposition}

The proof will be given elsewhere. The decompositions of that type were first found computationally: \href{https://www.kaggle.com/code/luxoove/lrx-optimal-algorithm-two-ways-bubble-sort}{notebook}.

%\subsection{The long elements,  and the formula for the longest one}


\subsection{Lower bound on the diameter}
Below, we provide a lower bound on the diameter using a combinatorial argument. The key idea is that cyclic shifts preserve the cyclic order of any given triple of numbers. Therefore, for each triple where the order is disrupted, the move $X$ is required in the decomposition.   The symmetries of a regular polygon, as described above, correspond precisely to the permutations that maximize changes in cyclic order. (The code to generate these long elements: \href{https://www.kaggle.com/code/alexandervc/lrx-long-elements-2-f-petrov-oeisa186752}{notebook}.) 

\begin{proposition}
    The diameter of the LRX Cayley graph of $S_n$ is larger than or equal to $n(n-1)/2-n/2-1$.
\end{proposition}

First, we reformulate the setup as follows.

The numbers $1,\ldots,n$  are written in the vertices of a regular $n$-gon $A_1A_2\ldots A_n$, each number being written once. At each move, one can either rotate the arrangement by $2\pi/n$  counterclockwise or clockwise, or swap the numbers at $A_1$ and $A_2$.
For two arrangements $\pi_1,\pi_2$ denote by $d(\pi_1,\pi_2)$ the minimal number of moves which suffice to get $\pi_2$ from $\pi_1$.

\textbf{Claim.} If $\pi_2$ is obtained from $\pi_1$ by an axial symmetry, then $d(\pi_1,\pi_2)\geqslant n^2/2-n-1$.

\textbf{Proof.} Assume that $\pi_2$ is obtained from $\pi_1$ by several moves. Call two numbers $a, b\in \{1,\ldots,n\}$ friends, if they were swapped odd number of times. Since for every three numbers $a, b, c$ the orientation of a triangle formed by these numbers changed, there is odd number of pairs of friends between $a, b, c$. Let $A$ denote the set of all non-friends of element 1, $B$ the complement  of $A$, i. e. $B$ consists of 1 and its friends. Our condition yields that all elements of $A$ are mutual friends, and so are all elements of $B$, but there is no friendship between $A$ and $B$. Thus, the total number of pairs of friends is $n(n-1)/2-|A|\cdot |B|\geqslant n(n-1)/2-n^2/4$. Hence, there were at least as many swaps. Between any two swaps there should be a rotation (otherwise, we are doing something useless). Totally, the number of operations is not less than $n(n-1)-n^2/2-1=n^2/2-n-1$.

\subsection{Upper bound on the diameter and the algorithms}
We develop two algorithms which can decompose any permutation in the product of LRX generators. We prove  complexity bound $n(n-1)/2+3n$, thus bounding diameter from above. The code and tests for the first algorithm can be found in \href{https://www.kaggle.com/code/mixnota/article-project}{notebook}.

Its brief description is the following:

The algorithm decomposes a permutation into the product of cycles. After that, it works with each cycle \( (a_1, a_2, \dots, a_n) \) in the following way:

1) We compute the permutation corresponding to the cycle.

2) We initialize the special variable \( x \).

3) The element \( a_1 \) is placed in the \( a_n \)-th position, and we perform a sequence of elementary transpositions to move \( a_1 \) from the \( a_n \)-th position to the \( a_1 \)-th position. During each elementary transposition, we either increase or decrease \( x \) by 1. The sign of the change depends on the direction of the elementary transposition. Essentially, \( x \) represents the current position of the transposition, but if our transposition is \( (1, N) \), the change in \( x \) is also \( \pm1 \), even though the position changes from 1 to \( N \) or from \( N \) to 1.

4) Next, we attempt to restore the position of \( a_2 \) and compute the sequence of elementary transpositions that move \( a_2 \) to the \( a_2 \)-th position. We apply the same operations to \( x \), and so on.

5) We repeat these actions for every element of the cycle.

6) In certain situations, the elementary transposition should be replaced by \( L \) or \( R \) (depending on the direction in which we are moving the current element). This replacement should occur if and only if the variable \( x \) satisfies the condition:

\[
x - c = \pm (N - 1)
\]


\textbf{Complexity estimation.}
First, when decomposing a cycle into transpositions, we do not shift anything by more than 1, except for elements already affected by the cycle.
Second, when moving an element from position \( i \) (or \( i \pm 1 \)) to position \( \sigma(i) \), we perform a sequence of at most \( \text{dist}(i, \sigma(i)) + 1 \) transpositions.
Third, the transition from \( i \to \sigma(i) \) to \( \sigma(i) \to \sigma(\sigma(i)) \) takes at most one rotation.
Thus, we perform at most  
\[
\sum 2(\text{dist}(i, \sigma(i)) + 1) + (\text{number of elements in the cycle})
\]
operations. The transitions between cycles and returning back will take at most \( n \) rotations.
Therefore, the entire decomposition takes no more than  
\[
2n + \sum 2(\text{dist}(i, \sigma(i)))
\]
operations.

Finally, after minimization,
\[
\sum \text{dist}(i, \sigma(i)) \leq \frac{n^2}{4}
\]
which gives a total of  
\[
2n + \frac{n^2}{2} + \frac{n}{2}
\]
operations.



We also develop a second algorithm: \href{https://www.kaggle.com/code/luxoove/top1-lrx-inversions-based-algorithm}{notebook}. Currently, it is our top performing algorithm from the practical tests, we expect its complexity is bounded by $n(n-1)/2+n/2$, however, it is not yet proved.

It might be expected that optimal algorithm with polynomial complexity can be developed, but it is not achieved yet. 

\subsection{Conjectures: Gumbel for growth, spectrum uniformity, random walks mixing, etc. }

Based on explicit computations for the growth for $n \le 15$ and its analysis (notebooks \href{https://www.kaggle.com/code/ogurtsov/gumbel}{1} \href{https://www.kaggle.com/code/ogurtsov/gumbel-for-binary-puzzle}{2} ) we come to the following conjecture, which can be thought as an analogue of the central limit theorems. It is in the vein of works by P.Diaconis et.al. ~\cite{diaconis1977spearman, diaconis1988metrics,chatterjee2017central} where for Coxeter generators growth and some other statistics were shown to follow Gaussian normal law in a limit. A new point of our conjecture is the appearance of asymmetric Gumbel distribution. Such asymmetry is quite natural since for the randomly chosen generators growth will be highly asymmetric with exponential growth and sharp abrupt (\href{https://mathoverflow.net/q/322877/10446}{Rubik's cube growth} is a prototypical example), while for nearly commutative Coxeter generators the growth is symmetric Gaussian. A slight asymmetry of growth for LRX generators indicates its intermediate nature between "exponential growth" and "commutative-like Gaussian growth" cases.

The growth distribution of the LRX graph seems to follow the Gumbel distribution for large $n$. The detailed analysis is provided in supplementary section~\ref{growth}.

\begin{figure}[ht!]
   \centering
   \includegraphics[width=1.0\linewidth]{figs/LRX_spec.png}
   \caption{Spectrum distribution of the LRX graphs. }\label{fig:LRX-spec}
\end{figure}

The analysis of the spectrum of LRX graphs is performed in \href{https://www.kaggle.com/code/nikolenkosergei/spectrum-analysis}{notebook}. From the figure above it is natural to expect that spectrum tends to uniform distribution.

We have computed all possible shortest paths from the conjecturally longest element to the identity by methods of the dynamical programming (\href{https://www.kaggle.com/code/fedmug/lrx-longest-paths}{notebook} . And after that analyzed averaged trajectories (\href{https://www.kaggle.com/code/antoninadolgorukova/lrx-sorting-networks}{notebook}) which are analogs of random sorting networks~\cite{angel2007random}. The figure \ref{fig:LRX-sort} below represents the results. It seems the pattern is different from the "sine curves" discovered in~\cite{angel2007random}. 

\begin{figure}[ht!]
   \centering
   \includegraphics[width=1.0\linewidth]{figs/LRX_sorting.jpg}
   \caption{Analogs of "random sorting networks" for LRX. Trajectories of individual elements under average of all possible shortest trajectories between longest node and identity. }\label{fig:LRX-sort}
\end{figure}

% We also numerically checked for even  $n \le 9$ (\href{https://www.kaggle.com/code/fedmug/lrx-shifts-starting-with-x}{notebook}) the following conjecture. It   would imply precise diameter lower bound $n(n-1)/2$. 

% {\bf Conjecture.} If the shortest  decomposition of an element starts from $X$, then its multiplications by $L^k$ $k<n/2$  have all different lengths.  

% Such a conjecture imply that half of the symmetries of the regular polygon have different lengths, but since their lengths greater than $n(n-1)/2-n/2$ one of the would have length $n(n-1)/2$ which is desired diameter. 

% P.Diacons and L.Soloff-Coste studied random walks on LRX Cayley~\cite{diaconis1993comparison}. They estimated that mixing time is between
% $n^2$ and $n^3\log(n)$. We performed numerical experiments to make more precise conjecture (notebooks \href{https://www.kaggle.com/competitions/lrx-oeis-a-186783-brainstorm-math-conjecture}{1},\href{https://www.kaggle.com/code/bbbxttt/wip-symmetric-groups-with-lrx-random-walk?scriptVersionId=220660303}{2}) unfortunately  the achieved range of $n$ was not enough to draw conclusions, however it seems that mixing time is more close to $n^3\log(n)$, moreover even for non-backtracking random walks, which is a bit surprising since they typically mix faster.

%Finally, we considered the Schreier coset graph for the action of $S_{2n}$ on binary strings with $n$ zeros and $n$ ones. (Analogue of $Gr(n,2n)$ over the "field with one element"). Efficient  \href{https://www.kaggle.com/code/ivankolt/lrx-4bit-uint64-coset}{code} which stores such strings in int64 allowed us to completely traverse the graph up to $n=42$ and compute its growth.
%The analysis of the data presented in the 
\href{https://www.kaggle.com/code/ogurtsov/gumbel-for-binary-puzzle}{notebook}. 
%Based on that it is tempting to conjecture that the diameters of such graphs are between $n^2/6$ and $n^2/5$, while growth again can be approximated by the Gumbel distribution.



\section{Code availability}

The code is available at the Kaggle platform, where it can be easily launched:
\href{https://www.kaggle.com/competitions/lrx-oeis-a-186783-brainstorm-math-conjecture/code}{LRX OEIS-A186783 brainstorm math conjecture}

The link above leads to one of the three public Kaggle challenges which we created to stimulate the research and interplay between artificial intelligence and mathematical communities, the other two:
\href{https://www.kaggle.com/competitions/lrx-binary-in-search-of-gods-number}{LRX discover math and God's algorithm}, \href{https://www.kaggle.com/competitions/lrx-discover-math-gods-algorithm}{LRX-binary: in search of God's number}. Kaggle infrastructure provides a convenient way to benchmark and compare different algorithms by just making the submissions to these challenges. The code, the data and discussions are stored at the same place. The code can be executed for free on Kaggle cloud servers. The first challenge is to decompose the conjecturally longest elements and to find their shortest decompositions. The participants have already achieved the length $n(n-1)/2$; thus, if we believe the conjecture is true, the optimum is found. The second challenge is to decompose random permutations and aims to find optimal algorithms (at the moment of writing optimum is not achieved). The third one asks the same but for the binary strings acted by LRX generators. 


% \section{Acknowledgments}
\section*{Acknowledgements}
A.C. is deeply grateful to M. Douglas for his interest in this work, engaging discussions, and invitation to present preliminary results at the Harvard CMSA program on "Mathematics and Machine Learning" in Fall 2024; to M. Gromov, S. Nechaev, and V. Rubtsov for their invitation to give a talk (\href{https://youtu.be/RkmBwlSyhfA?si=KgqQtRFaqx5ykd_s}{video}) at "Representations, Probability, and Beyond: A Journey into Anatoly Vershik World" at IHES, as well as for stimulating discussions. A.C. is grateful to J. Mitchel for involving into the Kaggle Santa 2023 challenge, from which this project originated, and to M.Kontsevich, Y.Soibelman,  S.Gukov, A. Hayat, T. Smirnova-Nagnibeda,  D.Osipov, V. Kleptsyn, G.Olshanskii, A.Ershler, J. Ellenberg, G. Williamson, A. Sutherland,  Y. Fregier, P.A. Melies, I. Vlassopoulos, F.Khafizov, A.Zinovyev, H.Isambert, T. Rokicki for the discussions, interest and comments, to his wife A.Chervova and daugther K.Chervova for support, understanding and help with computational experiments.  

We are deeply grateful to the many colleagues who have contributed to the CayleyPy project at various stages of its development, including: N.Bukhal, J.Naghiev, A.Lenin, E.Uryvanov,  A. Abramov, M.Urakov, A.Kuchin,  B.Bulatov,  F.Faizullin, A.Aparnev, O.Nikitina, A.Titarenko, U.Kniaziuk, D.Naumov, A.Krasnyi, S.Botman, A.Kostin,
R.Vinogradov, I.Gaiur, 
% added
% D.Gorodkov, A.Rozanov,V.Nelin,  A. Ogurtsov, A. Trepetsky, A. Dolgorukova, S. Lytkin, S. Ermilov, L. Grunvald, A. Eliseev, G. Annikov, M. Evseev, F. Petrov, N. Narynbaev, S. Nikolenko, S. Krymskii, R. Turtayev, S. Kovalev, N. Rokotyan, G. Verbyi, L. Shishina, A.Korolkova, D.Mamaeva,  S.Fironov, A.Lukyanenko,
K.Yakovlev, V.Shitov, E.Durymanov, A.Kostin, R.Magdiev, M.Krinitskiy, P.Snopov. 







\begin{thebibliography}{10}

% \bibitem{ChoquetChristodoulou}   Y. Choquet-Bruhat and D. Christodoulou, ``Existence of global solutions of the Yang-Mills, Higgs and spinor field equations in $3+1$ dimensions'', Annales scientifiques de l'\'{E}.N.S, 4 s\'erie, (1981), 481-500.

\bibitem[Douglas22]{douglas2022machine}
Douglas, Michael R. "Machine learning as a tool in theoretical science." Nature Reviews Physics 4.3 (2022): 145-146.

\bibitem[Charton19]{lample2019deep}
Lample, Guillaume, and François Charton. "Deep learning for symbolic mathematics." arXiv preprint arXiv:1912.01412 (2019).

\bibitem[Davies21]{davies2021advancing}
Davies, A., Veli{\v{c}}kovi{\'c}, P., Buesing, L., Blackwell, S., Zheng, D., Tomašev, N., ...  Kohli, P. (2021). Advancing mathematics by guiding human intuition with AI. Nature, 600(7887), 70-74.

\bibitem[Bao23]{bao2021polytopes}
Bao, J., He, Y. H., Hirst, E., Hofscheier, J., Kasprzyk, A., Majumder, S. (2023). Polytopes and machine learning. International Journal of Data Science in the Mathematical Sciences, 1(02), 181-211.


\bibitem[Romera24]{romera2024mathematical}
Romera-Paredes, B., Barekatain, M., Novikov, A., Balog, M., Kumar, M. P., Dupont, E., ...  Fawzi, A. (2024). Mathematical discoveries from program search with large language models. Nature, 625(7995), 468-475.

\bibitem[Coates23]{coates2024machine}
Coates, T., Kasprzyk, A.,  Veneziale, S. (2023). Machine learning detects terminal singularities. Advances in Neural Information Processing Systems, 36, 67183-67194.

\bibitem[Hayat25]{alfarano2024global}
Alfarano, A., Charton, F.,  Hayat, A. (2025). Global Lyapunov functions: a long-standing open problem in mathematics, with symbolic transformers. Advances in Neural Information Processing Systems, 37, 93643-93670.

\bibitem[Williamson24]{charton2024patternboost}
Charton, F., Ellenberg, J. S., Wagner, A. Z.,  Williamson, G. (2024). PatternBoost: Constructions in mathematics with a little help from AI. arXiv preprint arXiv:2411.00566.

\bibitem[Gukov24]{shehper2024makes}
Shehper, A., Medina-Mardones, A. M., Lewandowski, B., Gruen, A., Kucharski, P.,  Gukov, S. (2024). What makes math problems hard for reinforcement learning: a case study. arXiv preprint arXiv:2408.15332.

\bibitem[Kohli25]{swirszcz2025advancing}
Swirszcz, G., Wagner, A. Z., Williamson, G., Blackwell, S., Georgiev, B., Davies, A., ... Kohli, P. (2025). Advancing Geometry with AI: Multi-agent Generation of Polytopes. arXiv preprint arXiv:2502.05199.

\bibitem[AlphaGo16]{silver2016mastering}
Silver, D., Huang, A., Maddison, C. J., Guez, A., Sifre, L., Van Den Driessche, G., ...  Hassabis, D. (2016). Mastering the game of Go with deep neural networks and tree search. nature, 529(7587), 484-489.

\bibitem[AlphaZero17]{silver2017mastering}
Silver, D., Schrittwieser, J., Simonyan, K., Antonoglou, I., Huang, A., Guez, A., ... Hassabis, D. (2017). Mastering the game of go without human knowledge. nature, 550(7676), 354-359.

\bibitem[Chervov25]{chervov2025machinelearningapproachbeats}
Chervov, A., Khoruzhii, K., Bukhal, N., Naghiyev, J., Zamkovoy, V., Koltsov, I., ...  Romanov, A. (2025). A Machine Learning Approach That Beats Large Rubik's Cubes. arXiv preprint arXiv:2502.13266.

\bibitem[Babai89]{babai1989small}
Babai, L., Kantor, W. M.,  Lubotsky, A. (1989). Small-diameter Cayley graphs for finite simple groups. European Journal of Combinatorics, 10(6), 507-522.

\bibitem[Kuppili20]{kuppili2020upper}
Kuppili, S. S., Chitturi, B., Ravella, V. V.,  Datta, C. P. (2020, December). An Upper Bound for Sorting R n with LRE. In International Advanced Computing Conference (pp. 283-295). Singapore: Springer Singapore.

\bibitem[Diaconis77]{diaconis1977spearman}
Diaconis, P., Graham, R. L. (1977). Spearman's footrule as a measure of disarray. Journal of the Royal Statistical Society Series B: Statistical Methodology, 39(2), 262-268.

\bibitem[Diaconis88]{diaconis1988metrics}
Diaconis, P. (1988). Metrics on Groups, and Their Statistical Uses. In Group representations in probability and statistics (Vol. 11, pp. 102-131). Institute of Mathematical Statistics

\bibitem[Diaconis17]{chatterjee2017central}
Chatterjee, S.,  Diaconis, P. (2017). A central limit theorem for a new statistic on permutations. Indian Journal of Pure and Applied Mathematics, 48(4), 561-573.

\bibitem[Virág]{angel2007random}
Angel, O., Holroyd, A. E., Romik, D.,  Virág, B. (2007). Random sorting networks. Advances in Mathematics, 215(2), 839-868.

\bibitem[Diaconis93]{diaconis1993comparison}
Diaconis, P.,  Saloff-Coste, L. (1993). Comparison techniques for random walk on finite groups. The Annals of Probability, 2131-2156.

\bibitem[Goldreich81]{even1981minimum}
Even, S., Goldreich, O. (1981). The minimum-length generator sequence problem is NP-hard. Journal of Algorithms, 2(3), 311-313.

\bibitem[Jerrum85]{jerrum1985complexity}
Jerrum, M. R. (1985). The complexity of finding minimum-length generator sequences. Theoretical Computer Science, 36, 265-289.


\bibitem[Demaine17]{demaine2017solving}
Demaine, E. D., Eisenstat, S.,  Rudoy, M. (2017). Solving the Rubik's Cube Optimally is NP-complete. arXiv preprint arXiv:1706.06708.

\bibitem[Bulteau15]{bulteau2015pancake}
Bulteau, L., Fertin, G.,  Rusu, I. (2015). Pancake flipping is hard. Journal of Computer and System Sciences, 81(8), 1556-1574.

\bibitem[Veličković24]{wilson2024cayley}
Wilson, J. J., Bechler-Speicher, M., Veličković, P. (2024). Cayley graph propagation. arXiv preprint arXiv:2410.03424.

\bibitem[Korf97]{korf1997finding}
Korf, R. E. (1997, July). Finding optimal solutions to Rubik's Cube using pattern databases. In AAAI/IAAI (pp. 700-705).

\bibitem[Rokicki14]{Rokicki2014Diameter}
Rokicki, T., Kociemba, H., Davidson, M.,  Dethridge, J. (2014). The diameter of the rubik's cube group is twenty. siam REVIEW, 56(4), 645-670.

\bibitem[Sims70]{sims1970computational}
Sims, C. C. (1970, January). Computational methods in the study of permutation groups. In Computational problems in abstract algebra (pp. 169-183). Pergamon.

\bibitem[Knuth91]{knuth1991efficient}
Knuth, D. E. (1991). Efficient representation of perm groups. Combinatorica, 11(1), 33-43.

\bibitem[Shamir89]{fiat1989planning}
Fiat, A., Moses, S., Shamir, A., Shimshoni, I.,  Tardos, G. (1989, October). Planning and learning in permutation groups. In 30th Annual Symposium on Foundations of Computer Science (pp. 274-279). IEEE Computer Society.

\bibitem[Mulholland16]{mulholland2016permutation}
Mulholland, J. (2016). Permutation puzzles: a mathematical perspective. Departement Of mathematics Simon fraser University.

\bibitem[Babai88]{babai1988diameter}
Babai, L.,  Seress, Á. (1988). On the diameter of Cayley graphs of the symmetric group. Journal of combinatorial theory, Series A, 49(1), 175-179

\bibitem[Babai04]{babai2004diameter}
Babai, L., Beals, R.,  Seress, Á. (2004, January). On the diameter of the symmetric group: polynomial bounds. In SODA (pp. 1108-1112).


\bibitem[Seress14]{bamberg2014bounds}
Bamberg, J., Gill, N., Hayes, T. P., Helfgott, H. A., Seress, Á.,  Spiga, P. (2014). Bounds on the diameter of Cayley graphs of the symmetric group. Journal of Algebraic Combinatorics, 40(1), 1-22.

\bibitem[Helfgott14]{helfgott2014diameter}
Helfgott, H. A., Seress, Á. (2014). On the diameter of permutation groups. Annals of mathematics, 611-658.

\bibitem[Helfgott19]{helfgott2019growth}
Helfgott, H. A. (2019). Growth in linear algebraic groups and permutation groups: towards a unified perspective. Groups St Andrews 2017 in Birmingham, 455, 300.


\bibitem[Gates79]{gates1979bounds}
Gates, W. H., Papadimitriou, C. H. (1979). Bounds for sorting by prefix reversal. Discrete mathematics, 27(1), 47-57.

\bibitem[Pevzner95]{Pevzner1995human2mice}
Hannenhalli, S., Pevzner, P. A. (1995, October). Transforming men into mice (polynomial algorithm for genomic distance problem). In Proceedings of IEEE 36th annual foundations of computer science (pp. 581-592). IEEE.

\bibitem[Pevzner99]{Pevzner1999cabbage2turnip}
Hannenhalli, S.,  Pevzner, P. A. (1999). Transforming cabbage into turnip: polynomial algorithm for sorting signed permutations by reversals. Journal of the ACM (JACM), 46(1), 1-27.


\bibitem[Bulteau19]{bulteau2019parameterized}
Bulteau, L.,  Weller, M. (2019). Parameterized algorithms in bioinformatics: an overview. Algorithms, 12(12), 256.

\bibitem[Świta23]{swita2023solving}
Świta, R.,  Suszyński, Z. (2023). Solving Full N× N× N Rubik’s Supercube Using Genetic Algorithm. International Journal of Computer Games Technology, 2023(1), 2445335.

\bibitem[McAleer19]{mcaleer2019solving}
McAleer, S., Agostinelli, F., Shmakov, A. K.,  Baldi, P. (2019, January). Solving the rubik's cube with approximate policy iteration. In International Conference on Learning Representations.

\bibitem[Agostinelli19]{agostinelli2019solving}
Agostinelli, F., McAleer, S., Shmakov, A.,  Baldi, P. (2019). Solving the Rubik’s cube with deep reinforcement learning and search. Nature Machine Intelligence, 1(8), 356-363.

\bibitem[Agostinelli24]{khandelwal2024towards}
Khandelwal, V., Sheth, A.,  Agostinelli, F. (2024). Towards Learning Foundation Models for Heuristic Functions to Solve Pathfinding Problems. arXiv preprint arXiv:2406.02598.

\bibitem[Agostinelli24q]{agostinelli2024q}
Agostinelli, F., Shperberg, S. S., Shmakov, A., McAleer, S., Fox, R., Baldi, P. (2024). Q* search: Heuristic search with deep q-networks. In ICAPS Workshop on Bridging the Gap between AI Planning and Reinforcement Learning.

\bibitem[Takano21]{takano2023selfsupervision}
Takano, K. (2021). Self-Supervision is All You Need for Solving Rubik's Cube. arXiv preprint arXiv:2106.03157.

\bibitem[TakanBrunetto17]{brunetto2017deep}
Brunetto, R.,  Trunda, O. (2017, September). Deep Heuristic-learning in the Rubik's Cube Domain: An Experimental Evaluation. In ITAT (pp. 57-64).


\bibitem[Johnson21]{johnson2021solving}
Johnson, C. G. (2021). Solving the Rubik's cube with stepwise deep learning. Expert Systems, 38(3), e12665.

\bibitem[Amrutha22]{amrutha2022deep}
Amrutha, B. V.,  Srinath, R. (2022). Deep Learning Models for Rubik’s Cube with Entropy Modelling. In ICDSMLA 2020: Proceedings of the 2nd International Conference on Data Science, Machine Learning and Applications (pp. 35-43). Springer Singapore.


\bibitem[Noever21]{noever2022puzzle}
Noever, D.,  Burdick, R. (2021). Puzzle solving without search or human knowledge: An unnatural language approach. arXiv preprint arXiv:2109.02797.


\bibitem[Chasmai21]{chasmai2022cubetr}
Chasmai, M. E. (2021). CubeTR: Learning to Solve the Rubik's Cube using Transformers.

\bibitem[Bedaywi23]{bedaywi2023solving}
Bedaywi, Mark and Longtai, (Ted) Deng and Shaul, Francus (2023). Solving the Rubik’s Cube via Sequence Modeling using
Transformer-based Models, https://github.com/tedtedtedtedtedted/Solve-Rubiks-Cube-Via-Transformer/blob/main/Report.pdf



\bibitem[Pan21]{pan2021fourier}
Pan, H.,  Kondor, R. (2021, March). Fourier Bases for Solving Permutation Puzzles. In International Conference on Artificial Intelligence and Statistics (pp. 172-180). PMLR.

\bibitem[Swan17]{swan2017harmonic}
Swan, J. (2017, June). Harmonic analysis and resynthesis of Sliding-Tile Puzzle heuristics. In 2017 IEEE Congress on Evolutionary Computation (CEC) (pp. 516-524). IEEE.


\bibitem[Leskovec22]{pandy2022learning}
Pándy, M., Qiu, W., Corso, G., Veličković, P., Ying, Z., Leskovec, J.,  Liò, P. (2022, December). Learning graph search heuristics. In Learning on Graphs Conference (pp. 10-1). PMLR.

\bibitem[Yakovlev23]{kirilenko2023transpath}
Kirilenko, D., Andreychuk, A., Panov, A.,  Yakovlev, K. (2023, June). Transpath: Learning heuristics for grid-based pathfinding via transformers. In Proceedings of the AAAI Conference on Artificial Intelligence (Vol. 37, No. 10, pp. 12436-12443).

\bibitem[Gromov93]{gromov1993geometric}
Gromov, M. (1993). Geometric Group Theory: Asymptotic invariants of infinite groups (Vol. 2). Cambridge University Press.

\bibitem[Tao15]{tao2015expansion}
Tao, T. (2015). Expansion in finite simple groups of Lie type (Vol. 164). American Mathematical Soc..

\bibitem[Babai92]{babai1992diameter}
Babai, L.,  Seress, Á. (1992). On the diameter of permutation groups. European journal of combinatorics, 13(4), 231-243.

\bibitem[Akers89]{akers1989group}
Akers, S. B., Krishnamurthy, B. (1989). A group-theoretic model for symmetric interconnection networks. IEEE transactions on Computers, 38(4), 555-566.

\bibitem[Cooperman91]{cooperman1991applications}
Cooperman, G., Finkelstein, L.,  Sarawagi, N. (1991). Applications of cayley graphs. In Applied Algebra, Algebraic Algorithms and Error-Correcting Codes: 8th International Conference, AAECC-8 Tokyo, Japan, August 20–24, 1990 Proceedings 8 (pp. 367-378). Springer Berlin Heidelberg.

\bibitem[Heydemann97]{heydemann1997cayley}
Heydemann, M. C. (1997). Cayley graphs and interconnection networks. In Graph symmetry: algebraic methods and applications (pp. 167-224). Dordrecht: Springer Netherlands

\bibitem[Wigderson06]{hoory2006expander}
Hoory, S., Linial, N.,  Wigderson, A. (2006). Expander graphs and their applications. Bulletin of the American Mathematical Society, 43(4), 439-561.

\bibitem[Zémor94]{zemor1994hash}
Zémor, G. (1994). Hash functions and Cayley graphs. Designs, Codes and Cryptography, 4(3), 381-394.

\bibitem[Petit11]{petit2013rubik}
Petit, C., Quisquater, J. J. (2011). Rubik's for cryptographers. Cryptology ePrint Archive.

\bibitem[Kohli24]{ruiz2024quantum}
Ruiz, F. J., Laakkonen, T., Bausch, J., Balog, M., Barekatain, M., Heras, F. J., ...  Kohli, P. (2024). Quantum Circuit Optimization with AlphaTensor.(2024). URL https://arxiv. org/abs/2402.14396.

\bibitem[Sarkar24]{sarkar2024quantum}
Sarkar, R. S.,  Adhikari, B. (2024). Quantum circuit model for discrete-time three-state quantum walks on Cayley graphs. Physical Review A, 110(1), 012617.


\bibitem[Vidick23]{dinur2023good}
Dinur, I., Hsieh, M. H., Lin, T. C.,  Vidick, T. (2023, June). Good quantum LDPC codes with linear time decoders. In Proceedings of the 55th annual ACM symposium on theory of computing (pp. 905-918).

\bibitem[Acevedo06]{acevedo2006exploring}
Acevedo, O. L., Roland, J.,  Cerf, N. J. (2006). Exploring scalar quantum walks on Cayley graphs. arXiv preprint quant-ph/0609234.

\bibitem[Gromada22]{gromada2022some}
Gromada, D. (2022). Some examples of quantum graphs. Letters in Mathematical Physics, 112(6), 122.

\bibitem[Helfgott15]{helfgott2015random}
Helfgott, H. A., Seress, Á.,  Zuk, A. (2015). Random generators of the symmetric group: diameter, mixing time and spectral gap. Journal of Algebra, 421, 349-368.


\bibitem[Diaconis13]{diaconis2013some}
Diaconis, P. (2013). Some things we’ve learned (about Markov chain Monte Carlo).

\bibitem[Hirata24]{hirata2024probabilistic}
Hirata, S. (2024). Probabilistic estimates of the diameters of the Rubik's Cube groups. arXiv preprint arXiv:2404.07337.

\bibitem[Cormen09]{cormen09} Cormen, T., Leiserson, Ch., Rivest, R., Stein, Cl. (2009). Introduction to Algorithms. MIT Press.

\bibitem[Sutton18]{sutton2018} Sutton, R., Barto, A. (2018). Reinforcement Learning: An Introduction. MIT Press.

\bibitem[Mnih13]{mnih2013playing} Mnih, V., Kavukcuoglu, K., Silver, D., Graves, A., Antonoglou, I., Wierstra, D., Riedmiller, M. (2013). Playing Atari with Deep Reinforcement Learning. \href{https://arxiv.org/abs/1312.5602}{ArXiv preprint}.

\bibitem[Alon07]{alon2007non}
Alon, N., Benjamini, I., Lubetzky, E.,  Sodin, S. (2007). Non-backtracking random walks mix faster. Communications in Contemporary Mathematics, 9(04), 585-603.

\bibitem[MacSturm]{MacSturm} Diane Maclagan and Bernd Sturmfels. \textit{Introduction to Tropical Geometry.} AMS, 2015. ISBN: 978-1-4704-6856-9. URL: https://bookstore.ams.org/view?ProductCode=GSM/161

\bibitem[Glushkov68]{glushkov1968completeness}
Glushkov, V. M. (1968). Completeness of a system of operations in digital computers. Cybernetics and Systems Analysis, 4(2), 1–5.
\href{https://link.springer.com/article/10.1007/BF01073731}{Springer}

\bibitem[GlukhovZubov99]{glukhovzubov1999lengths}
Glukhov, M. M., Zubov, A.Yu. (1999). “On the lengths of symmetric and alternating permutation groups in various generating systems (Review)". *Mathametical questions on cybernetics*, 8, 5–32. (In Russian).
\href{https://keldysh.ru/papers/1999/mvk/mvk1999_5.pdf}{Link Keldysh}, 
\href{https://www.kaggle.com/datasets/alexandervc/cayleypy-development-3-growth-computations/data?select=Glukhov_Zubov1999review.pdf}{Link Kaggle}

\bibitem[NijenhuisWilf75]{NijenhuisWilf1975CombinatorialAlgorithms}
Nijenhuis, A., Wilf, H. S. (1975).
Combinatorial Algorithms.
Academic Press, New York.

\bibitem[SawadaWilliams19]{SawadaWilliams2019SigmaTau}
Sawada, J., Williams, A. (2019).
Solving the Sigma-Tau Problem.
ACM Transactions on Algorithms, 16(1), Article 1-17.
doi:10.1145/3359589

\bibitem[Douglas25]{douglas2025diffusion}
Douglas, M. R., Fraser-Taliente, K. (2025). Diffusion Models for Cayley Graphs. *arXiv preprint* arXiv:2503.05558.

\bibitem[Ziarko25]{ZiarkoBortkiewiczZawalskiEysenbachMilos2025CRTR}
Ziarko, A., Bortkiewicz, M., Zawalski, M., Eysenbach, B., Milos, P. (2025).
Contrastive Representations for Temporal Reasoning.
arXiv:2508.13113.

\bibitem[HashemiCorominasGiacchetto25]{hashemi2025transformers}
Hashemi, B., Corominas, R. G., Giacchetto, A. (2025). *Can Transformers Do Enumerative Geometry?* arXiv preprint arXiv:2408.14915.

\bibitem[CayleyPyGrowth]{CayleyPyGrowth}
A. Chervov, D. Fedoriaka, E.V. Konstantinova, A. Naumov, I. Kiselev, A. Sheveleva, I. Koltsov, S. Lytkin, A. Smolensky, A. Soibelman, F. Levkovich-Maslyuk, R. Grimov, D. Volovich, A. Isakov, A. Kostin, M. Litvinov, N. Vilkin-Krom, A. Bidzhiev, A. Krasnyi, M. Evseev, E. Geraseva, L. Grunwald, S. Galkin, E. Koldunov, S. Diner, A. Chevychelov, E. Kudasheva, A. Sychev, A. Kravchenko, Z. Kogan, A. Natyrova, L. Shishina, L. Cheldieva, V. Zamkovoy, D. Kovalenko, O. Papulov, S. Kudashev, D. Shiltsov, R. Turtayev, O. Nikitina, D. Mamayeva, S. Nikolenko, M. Obozov, A. Titarenko, A. Dolgorukova, A. Aparnev, O. Debeaupuis, S. Alami C., H. Isambert,
\textit{CayleyPy Growth: Efficient growth computations and hundreds of new conjectures on Cayley graphs (Brief version).}
arXiv preprint arXiv:2509.19162, 2025.
URL: \url{https://arxiv.org/abs/2509.19162}.

\end{thebibliography}

\section{Supplementary}

Here we present additional analysis related both to mathematical and AI counterparts presented in the main text.
% numerical experiments exploring various properties of LRX Cayley graphs. 
% Based on these results, we propose several mathematical conjectures.
Subsection \ref{growth} we give details on the growth approximation by Gumbel distribution, the next subsection devoted to similar analysis
for the Schreier coset graph and formulating conjectures for it.
Subsection \ref{sec-mix} we analyze X-trick and random walks, in particular we conjecture that mixing time for random walks with X-trick is $O(n^2)$,
that means X-trick speed-ups mixing comparing to standard random walks which mixing time is known to be between $O(n^3)$ and $ O(n^3log(n))$. 
Which partially explains why X-trick is so effective in beam search.
Subsection \ref{sec-dif-dist} gives details on the diffusion distance and provides results confirming that it is not fully perfect heuristics, which justifies our main algorithmic proposal - combining it with RL methods.
Subsection \ref{sec-beamsize-and-depth} presents analysis how performance depends on the beam width and depth of non-backtracking. 
(More simulation results for various parameters can be found in \href{https://www.kaggle.com/competitions/lrx-oeis-a-186783-brainstorm-math-conjecture/discussion/612338}{spreadsheet here}, or  on \href{https://docs.google.com/spreadsheets/d/1KFYTqPDSdH8_vdvWz3YLTLKg5IlW1Kkx_9kJClSKGOs/edit?usp=sharing}{google sheets}.)
In subsection \ref{sec-tropical} we discuss relation of the tropical mathematics and the Bellman equation.  

\subsection{Gumbel growth distribution}
\label{growth}

A \textit{layer} of a Cayley graph is the set of all vertices with the same distance to the target (e.g. identity permutation $e$). LRX Cayley graph for $S_n$ conjecturally consists of $\frac{n(n-1)}2 + 1$ layers. Using the BFS algorithm, we calculate the sizes of all these layers for $n\le 15$. Being normalized, the layer sizes form the \textit{growth distribution} of the graph. We hypothesize that for $n \to \infty$ it converges to the left-skewed Gumbel distribution with pdf:

\begin{equation}
\label{eq:gumbel}
    p(x \vert \mu, \beta) = \frac 1\beta F\Big(\frac{x-\mu}\beta \Big), \quad F(t) = e^{t - e^t}.
\end{equation}

\begin{figure}[ht!]
   \hspace*{-13mm}
   \includegraphics[width=1.2\linewidth]{figs/growth.png}
   \caption{Growth distribution of LRX Cayley graph fitted by left-skewed Gumbel~\eqref{eq:gumbel}.}\label{fig:LRX-growth}
\end{figure}

\begin{figure}[ht!]
   %\centering
   \hspace*{-15mm}
   \includegraphics[width=1.2\linewidth]{figs/LRX_stats.png}
   \caption{LRX growth statistics. Diameter, mean, mode, $\mathrm{std}^{3/2}$ seem to be quadratic polynomials of $n$. Skewness is a slightly negative monotonic function of $n$.}\label{fig:LRX-stats}
\end{figure}

To estimate the parameters $\mu$ and $\beta$, we calculate the mean and variance of the LRX growth distribution and equal them to $\mu - \gamma \beta$ and $\frac{\pi^2}6 \beta^2$, respectively (theoretical moments of~\eqref{eq:gumbel}). After that, we search by grid for slightly better parameters. The approximation obtained for the growth of LRX by the Gumbel distribution is presented in Fig.~\ref{fig:LRX-growth}. 

The mean, mode and standard deviation raised to the power of $\frac 32$ seem to be quadratic functions of $n$ (see Fig.~\ref{fig:LRX-stats}). The comparison of these LRX statistics with those obtained by the Gumbel fit is presented in Table~\ref{tab:growth_stats}. In all cases we observe a good fit by quadratic polynomials with higher coefficients within $[0.1, 0.5]$. See \href{https://www.kaggle.com/code/fedmug/lrx-growth-distribution}{notebook} for details.


The skewness and excess kurtosis are shown in the right part of Fig.~\ref{fig:LRX-stats}. Although they are not very close to the theoretical values of the Gumbel distribution, the skewness, not converging to $0$, reveals that the limit growth distribution is not Gaussian. 

% If the conjecture for the standard deviation is true, then $\beta^{3/2} \sim Cn^2$, $n\to\infty$, $\beta \sim C{n^{4/3}}$, and the maximum of the density~\eqref{eq:gumbel} is $p(\mu) = \frac 1{\beta e} \sim \frac 1{Ce n\sqrt[3]n}$. Hence, the richest layer of the LRX graph should contain about $Ce\frac{(n-1)!}{\sqrt[3] n}$ elements.
%The mode $\mu$ seems to be a quadratic function of $n$, $\mu(n) \approx 0.385 n^2$, which is approximately $77\%$ of the diameter. We also observed that $\beta(n) \sim \frac{n\sqrt n}{Ce}$ where $C$ is around $3$. Since maximum of the density~\eqref{eq:gumbel} is $p(\mu) = \frac 1{\beta e}$, from this assumption we conclude that the richest layer of the LRX graph should contain about $Ce\frac{n!}{n^{3/2}}$ elements.


\begin{table}[ht!]
    \centering
    \begin{tabular}{|c|l|l|}
        \hline
        \textbf{Statistic} & \textbf{LRX growth} & \textbf{Gumbel fit} \\
        \hline
        diameter & $0.5n^2 - 0.5n$ & ---\\
        mean & $0.3763n^2 - 1.0698n + 1.3015$ & $0.3765n^2 - 1.1205n + 1.3647$\\
        mode & $0.3896 n^2 - 0.983n + 1.0335$ & $0.3858 n^2 - 0.9545n + 1.1767$ \\
        $\mathrm{sd}^{3/2}$ & $0.1455n^2 - 0.7245n + 2.4859$ & $0.1686 n^2 - 0.6196n + 1.6563$ \\
        \hline
    \end{tabular}
    \caption{Quadratic fits for statistics of LRX growth.}
    \label{tab:growth_stats}
\end{table}

\clearpage
\subsection{Coset graph}
A similar analysis is performed for the Schreier coset graph for $S_{n}, n=2m$ on binary strings with $m$ zeros and $m$ ones (analogue of $Gr(m,2m)$ over the "field with one element"), $4\le m \le 21$.  Efficient \href{https://www.kaggle.com/code/ivankolt/lrx-4bit-uint64-coset}{code} that stores such strings in \texttt{int64} allowed us to traverse the graph up to $n=42$ and compute its growth. Once again, 
 it can be well approximated by the Gumbel distribution~\eqref{eq:gumbel} (Fig.~\ref{fig:coset-growth}).
Note: in case of Schreier graphs growth is dependent on the starting point in general - we take it to be first $m$ zeros, than $m$ ones.
In contrast to Cayley graphs we will speak not about the diameters of these graphs, but the about distance to the most far nodes from that
starting point - we will call them God's numbers - using terminology adopted in puzzle community.
% The most far elements from it - correspond to

\begin{figure}[ht!]
   \hspace*{-13mm}
   \includegraphics[width=1.2\linewidth]{figs/coset_growth.png}
   \caption{Growth distribution of the coset graph fitted by left-skewed Gumbel~\eqref{eq:gumbel}.}\label{fig:coset-growth}
\end{figure}

{\bf Conjectures.}
The growth  coset growth converges to the Gumbel~\eqref{eq:gumbel} in distribution for $n\to \infty$. 
The coset growth statistics for mean, mode can be approximately described by quadratic as presented in Fig.~\ref{fig:coset-stats} and Table.~\ref{tab:coset_growth_stats}, respectively.
Moreover it is tempting to conjecture that the God's numbers $d_n$ is (which surprisingly coincide with 
\href{https://oeis.org/A132297}{A132297}( Number of distinct Markov type classes of order 2 possible in binary strings of length n)
and with \href{https://oeis.org/A267529}{A267529}(Total number of ON (black) cells after n iterations of the "Rule 141" elementary cellular automaton starting with a single ON (black) cell)\footnote{We thank Tomas Rokicki for that observations}: 
$$
    d_n = \frac 3{16}n^2 - \frac n4 + 2 - \frac{n \mod 4}8, n \geq 6, n \text{ is even},
$$


See  (\href{https://www.kaggle.com/code/fedmug/lrx-coset-growth}{notebook}) for more details.
(God's number is quasi-polynomial - which is consistent which general conjecture presented in the subsequent paper:
\cite{CayleyPyGrowth}).
%the diameters of such graphs are between $n^2/6$ and $n^2/5$, 
% \noindent

\begin{figure}[ht!]
   %\centering
   \hspace*{-15mm}
   \includegraphics[width=1.2\linewidth]{figs/coset_stats.png}
   \caption{Coset growth statistics.}\label{fig:coset-stats}
\end{figure}

 \begin{table}[ht!]
    \centering
    \begin{tabular}{|c|l|l|}
        \hline
        \textbf{Statistic} & \textbf{Coset growth} & \textbf{Gumbel fit} \\
        \hline
        God's number & $0.1875n^2 - 0.25n + O(1)$ & ---\\
        mean & $0.1591 n^2 - 0.8887n + 3.4809$ & $0.1595n^2 - 0.9919n + 3.9014$\\
        mode & $0.169n^2 - 0.936n + 3.8898$ & $0.1657n^2 - 0.7831n + 3.2604$ \\
        $\mathrm{sd}^{3/2}$ & $0.1597 n^2 - 1.8523n + 8.3784$ & $0.1927n^2 - 1.757n + 6.6497$ \\
        \hline
    \end{tabular}
    \caption{Quadratic fits for statistics of the Schreier coset growth.}
    \label{tab:coset_growth_stats}
\end{table}

{\bf Long elements.} Currently we are unable to find a clear pattern for the longest elements, 
however we are able to identify some elements which are quite long.
The sizes of the last layer of the coset graph, which contain the longest elements, seem to have no clear pattern (Table~\ref{tab:coset_longest}). The projection of the conjecturally longest element~\eqref{eq:longest} of the LRX Cayley graph onto the binary coset by the rule $ \big[\ell_n^i < \frac n2\big]$, $0 \leq i <n$, produces
$00\underbrace{11\ldots 1}_{\lceil \frac n2 \rceil}\underbrace{00\ldots 0}_{\lfloor \frac n2 \rfloor - 2}
$,
whose length obviously does not exceed $\frac n2 - 2$ and hence these elements (longest in the full Cayley graph) cannot be longest in the coset graph. We propose a family of long elements whose length for feasible $n$ is close to the God's numbers $d_n$:

\begin{equation}
    c_n = \underbrace{11\ldots 1}_{n_0}\underbrace{00\ldots 0}_{n_1} \underbrace{11\ldots 1}_{n_2}\underbrace{00\ldots 0}_{n_3}, \quad n_i = \Big\lfloor \frac{n+i}4\Big\rfloor.
\end{equation}

We examine these elements and all their cyclic shifts, and calculate the distance to the central element $00\ldots 011\ldots 1$. The maxima of these distances, presented in the last column of Table~\ref{tab:coset_longest}, turn out to be close to the diameter $d_n$.

{\bf Conjecture.} The lengths of these elements are not smaller than God's number minus $O(n)$.  

\begin{table}[ht!]
\centering
\begin{tabular}{|c|c|c|c|c|c|}
\hline
$n$ & \#Longest & God's numbers $d_n$ & $\mathrm{len}(c_n)$ & $\max\limits_{0\leq k < n}\mathrm{len}(L^kc_n)$ & $\max\limits_{s \leftrightsquigarrow c_n} \mathrm{len}(s)$ \\
\hline\hline
4  & 3  & 2  & 2 & 2 & 2 \\ \hline
6  & 1  & 7  & 3 & 6 & 7  \\ \hline
8  & 1  & 12 & 10 & 10 & 10 \\ \hline
10 & 4  & 18 & 13 & 17 & 18 \\ \hline
12 & 4  & 26 & 23 & 24 & 24 \\ \hline
14 & 11 & 35 & 27 & 32 & 34 \\ \hline
16 & 6  & 46 & 42 & 43 & 43 \\ \hline
18 & 14 & 58 & 47 & 53 & 57 \\ \hline
20 & 10 & 72 & 67 & 69 & 69 \\ \hline
22 & 32 & 87 & 73 & 80 & 84 \\ \hline
24 & 16 & 104 & 98 & 100 & 100 \\ \hline
26 & 49 & 122 & 105 & 113 & 119 \\ \hline
28 & 26 & 142 & 135 & 138 & 138 \\ \hline
30 & 99 & 163 & 141 & 152 & 158 \\ \hline
32 & 42 & 186 & 178 & 181 & 181 \\
\hline
\end{tabular}

\caption{Number of longest elements, the length of the long (but not the longest) elements $c_n$, maximal length of their shifts and the length of the longest element reachable from these shifts by distance non decreasing steps  for coset graphs. The values in the last three columns seem to be as close to $d_n$ as $O(n)$.}
\label{tab:coset_longest}
\end{table}

\clearpage
\subsection{Mixing time}\label{sec-mix}

P.Diacons and L.Soloff-Coste studied random walks on LRX Cayley graphs~\cite{diaconis1993comparison}. 
They estimated the \href{https://en.wikipedia.org/wiki/Markov_chain_mixing_time}{\textit{mixing time}} --- the number of steps needed for a permutation to resemble a uniform random sample --- to be between
$n^3$ and $n^3\log n$.
Here we complement their fundamental results by computational analysis of the different types of the random walks on the same LRX Cayley graph
and formulate new conjectures (items 2,3): 

\begin{enumerate}
\item 
pure random walks (all moves L, R, X are equiprobable on each step) - mixing time is between $n^3$ and $n^3\log n$~\cite{diaconis1993comparison}
\item 
random walks with non-backtracking (forbids going back where it just came from) - conjecture: mixing time has the same order  between $n^3$ and $n^3\log n$
\item 
random walks with non-backtracking and X-trick -  conjecture: mixing time reduces to  $n^2$ 
\end{enumerate}

In order to analyze the behavior of the random walks we consider the following statistics - number of inversions.
The advantage of that statistics that it can be computed very fast. 
Figure \ref{fig:LRX-mixing} presents values of that statistics averaged over  $5000$ random walks.
% In our experiments we run three series of $5000$ random walks from the identity permutation:
Roughly speaking the time (step) of achieving the plateau corresponds to the mixing time,
more precisely such time is lower bound for the mixing time. 
(As expected for the first two series the level of the plateau is $\frac{n(n-1)}4$ which coincides with expectation of the number of inversions of a random permutation, while for X-tricks random walks it is different). 

To estimate the mixing time we used the step when a series of random walks reaches 95\% or 99\% of the asymptote (\href{https://www.kaggle.com/code/fedmug/lrx-random-walk-asymptote}{notebook}). For random walks of type~(1) and~(2) these time steps are fitted by a cubic polynomial - that corresponds to mixing time being   between $n^3$ and $n^3\log n$.  Interestingly, the best fit for the series~(3) is delivered by a quadratic polynomial. Consequently, our conjecture for the mixing time is $O(n^3)$; introducing X-trick allows to mix faster in $O(n^2)$, however, the distribution in this case is clearly not uniform. Such speed-up of mixing may serve as an explanation why X-trick is so effective in beam-search.


\begin{figure}[ht!]
   \centering
   \includegraphics[width=1.0\linewidth]{figs/inversions_avg.png}
   \caption{Pure random walks, random walks with non-backtracking, and random walks with non-backtracking and X-trick. The latter mixes faster than the first two.}\label{fig:LRX-mixing}
\end{figure}

% Fig.~\ref{fig:LRX-mixing} shows the average number of inversions as function of step in each series. In all cases this function has a horizontal asymptote. For the first two series the level of the plateau is $\frac{n(n-1)}4$ which coincides with expectation of the number of inversions of a random permutation.


%We performed numerical experiments to make more precise conjecture (notebooks \href{https://www.kaggle.com/competitions/lrx-oeis-a-186783-brainstorm-math-conjecture}{1},\href{https://www.kaggle.com/code/bbbxttt/wip-symmetric-groups-with-lrx-random-walk?scriptVersionId=220660303}{2}) unfortunately  the achieved range of $n$ was not enough to draw conclusions, however it seems that mixing time is more close to $n^3\log(n)$, moreover even for non-backtracking random walks, which is a bit surprising since they typically mix faster.

\clearpage
\subsection{LRX diffusion distance analysis}\label{sec-dif-dist}

Diffusion distance is a very important part of our AI framework. 
Roughly speaking it should be thought as a length of a random path between two nodes. 
That phrasing is almost correct, but hides the  detail which is quite important. There is actually no "THE" diffusion distance,
there is no way to define one canonical diffusion distance on a graph, it is always a family depending on some parameter. 
A conceptual explanation: 
in a sense diffusion distance is a deformation of the canonical distance (length of the shortest path),
but since it is a deformation it must depend on some parameter. 
The most straight forward way to introduce the parameter is the following:
we fix some number of steps $K_{max}$ and consider random paths of the length $K_{max}$, then the average number of steps,
which such paths needs to get from node $A$ to node $B$ corresponds to the diffusion distance. 
Such way is similar to \href{https://en.wikipedia.org/wiki/Brownian_bridge}{Brownian bridge}  -  one needs to fix time $T$ to define it - that is similar to our fixation of the $K_{max}$.  

{\bf Definition of the diffusion distance. } Consider a graph, let $S$ be some selected node (e.g. "solved state", identity in group). 

Let us denote by $P(V,K)$ the transition probability from node $S$ to node $V$ in $K$ steps. I.e.: 
\begin{equation}
    P(V,K) = \text{Probability to reach V from S in K steps }
    \label{eq:P}
\end{equation}
That is -- one considers all possible paths of length $K$ and starting from $S$. And one divides number of such paths ending at $V$, by the number all such paths. (In particular we may get zero if $V$ is not reachable it $K$ steps.)

 Diffusion distance to node $V$ from node $S$ for the random walks length fixed to be  $K_{max}$ can be defined:
\begin{equation}
    \mathrm{DD}(V,K_{max}) = \frac{\sum_{k=0,...,K_{max} } ~  k ~  P(V,k)  }{\sum_{k=0,...,K_{max} }  P(V,k)}
\end{equation}
The sense of the formula above is clear -  it is weighted average of numbers of steps of random walks: $1,2,3,...K_{max}$ where weights are given by probabilities to reach $V$ in these numbers of steps. So it is exactly "average length of a random path from S to V" in certain precise sense and dependence on $K_{max}$ is explicit. The numerator and denominator are different only by presence  of $k$ in numerator. 



{\bf Formula for the diffusion distance. } (Case of simple random works - without non-backtracking or any other modification). 
Diffusion distance %(at least for the case of simple random walks) 
can be computed via the adjacency matrix of a graph as follows: 

\begin{equation}
    \label{eq:dd}
    %\mathrm{DD}(V,K_{max} ) = \frac{\sum\limits_{k=0}^{K_{\max}} k \mathbb P(V, k)}{\sum\limits_{k=0}^{K_{\max}}  \mathbb P(V, k)}
    \mathrm{DD}(V,K_{max} ) = \frac{\sum\limits_{k=0}^{K_{\max}} k  P(V, k)}{\sum\limits_{k=0}^{K_{\max}}   P(V, k)}
\end{equation}
    
\noindent
where $\mathbb P(v, k)$ is the probability to reach $v$ from the initial state $e$ in $k$ steps. If $A$ is the adjacency matrix of the Cayley graph then $\mathbb P(v, k) = \frac{A^k_{0v}}{\sum\limits_j A^k_{0j}}$. 
The formula reveals a striking similarity of diffusion distance to zeta functions, and supersymmetric Landau-Ginzburg models. We will discuss it in subsequent publications. 

{\bf Computations and analysis. } 
Calculation of powers of the adjacency matrix could be infeasible even for relatively small $n$. To speed up the computations, we use dynamic programming to calculate $A^k_{0v}$ --- number of paths of length $k$ from the identity permutation $e$ to the state $v$.
Results of the computations presented on the figure \ref{fig:LRX-dd-10}.
Conclusions: first diffusion distance has a huge variation when the actual distance (length of the shortest path is fixed),
one can see that histograms of the diffusion distance for layers $k$, $k+1$ intersect - which is unpleasant feature for diffusion
distance utilization in path-finding, since it will mislead  not to approach the goal, but deviate from it in some cases.
Second moreover, in some cases the same happens  not only for exceptional cases, but even in average diffusion distance 
for e.g. layer $22$ is smaller than for the layer $21$. 
These are the features demonstrating that diffusion distance is a not a perfect predictor of the true distance,
however from practical experience we see that it is working quite well in many cases, but still RL approach presented in the present paper
might be more promising. 


\begin{figure}[ht!]
\centering
\includegraphics[width=1.0\linewidth]{figs/dd_10_hist.png}
\caption{Histogram of diffusion distances for layers $20$--$23$ of LRX Cayley graph for $S_{10}$. The average diffusion distance for layer $22$ is clearly less than for layer $21$ - demonstrating non-monotonicity of the diffusion distance with respect to the actual distance. }\label{fig:LRX-dd-10}
\end{figure}

% For LRX Cayley graphs, the diffusion distance turned out to be a non-monotonic function of the true distance (see fig.~\ref{fig:LRX-dd-10}). Such behavior, discovered for all $n\le 11$ (\href{https://www.kaggle.com/code/fedmug/lrx-diffusion-distance-analysis}{notebook}), is a possible reason why our methods often could not find any path to the target in LRX Cayley graph beyond $n=32$.

\subsection{Impact of beam width and non-backtracking on the beam search performance}\label{sec-beamsize-and-depth}

The key parameter of beam search is  beam width (beam size). But advanced versions introduce various other parameters,
in particular non-backtracking depth, which controls how many previous moves are forbidden from being immediately revisited (to avoid short cycles). Varying these parameters affects performance: both the success rate and the length of the solution paths. Here describe results of such analysis. 

To study how beam width and non-backtracking depth affect the performance of ML-guided beam search, we conducted over 15,000 experiments (30–50 runs per condition) using XGBoost, the fastest model to date. We varied the permutation size ($n$), the beam width, and the non-backtracking depth (\href{https://www.kaggle.com/code/antoninadolgorukova/lrx-rw-bs-parameters-in-gbdt-guided-beam-search}{kaggle notebook}).
We tested all approaches on the element of the conjecturally longest element of length $n(n-1)/2$ described in previous sections.
These experiments revealed that increasing both parameters consistently improves the success rate, though gains diminish for larger $n$. 

To isolate the effect of non-backtracking, we set the beam width to the minimum value that yielded a non-zero success rate for each $n$ (15, 18, 20, 25, and 28) and systematically increased the depth of non-backtracking. As shown in Fig. \ref{fig:nbt_xgboost}A, increasing the depth of non-backtracking (expressed as a percentage of the conjectured optimal solution length, $n(n-1)/2)$) improves the success rate, although the magnitude of the improvement is smaller for larger permutations. In addition, deeper non-backtracking typically results in longer median solution lengths (Fig. \ref{fig:nbt_xgboost}B).
However, one cannot directly interpret that phenomena as decreasing performance, 
since success rate is improving - that beam search is able to solve cases more difficult for it,
that difficulty  manifest itself in  larger lengths and that might be the reason why  median solution length is growing. 


\begin{figure}[ht!]
  \centering
  \includegraphics[width=\textwidth]{figs/nbt_xgboost_2panel.png}
  \caption{Effect of non-backtracking (NBT) depth for beam search on (A) success rate and (B) median solution length.
  Beam width is fixed at the minimum viable value for each permutation size $n$.}
  \label{fig:nbt_xgboost}
\end{figure}

\begin{figure}[ht!]
  \centering
  \includegraphics[width=\textwidth]{figs/bw_xgboost_2panel.png}
  \caption{Effect of beam width on (A) success rate and (B) median solution length.
  Non-backtracking depth is fixed at two steps. Beam-width index $(i)$ is used
  instead of absolute beam widths: $i=1$ denotes the smallest beam width tested
  for each $n$, and each increment multiplies the beam width by two
  ($B_i = B_1 \cdot 2^{\,i-1}$).}
  \label{fig:bw_xgboost}
\end{figure}

To illustrate the impact of beam width on success rates, we used different beam widths at a fixed non-backtracking depth of two steps. Note that at step \(s\), we expand \(L_{s-1}\) to 1-move neighbors and discard any candidate whose hash appears in the stored layers (the last \(D\) frontier layers, where \(D\) is the non-backtracking depth). Since a 1-move neighbor
can only return to the immediate parent in \(L_{s-2}\), a 2-step non-backtracking depth (i.e., \(D=2\)) blocks exactly parent reversals; a 1-step depth (\(D=1\)) does not, because \(L_{s-2}\) lies outside the history. For each permutation size $n$ (15, 18, 20, 25, and 28), we started from the minimum beam width that yielded a non-zero success rate. Then, we increased it exponentially (in powers of two) 5 to 6 times. Since the minimal successful beam widths differ dramatically between small and large $n$ (e.g., 2 for $n = 15$ and 16384 for $n = 28$), we represented beam widths by their increment index $i$ rather than absolute values. Thus, $i=1$ always denotes the smallest beam width tested for each $n$, $i=2$ denotes the second smallest, and so forth. This indexing allowed us to plot success rates for all permutation sizes on a unified x-axis (Fig. \ref{fig:bw_xgboost}A).

This experiment clearly demonstrates a consistent trend within each $n$: increasing the beam width leads to higher success rates. However, the marginal gains diminish as $n$ increases. That is, larger permutations require significantly wider beams to achieve the same improvement in the success rate, indicating reduced search efficiency in higher-dimensional search spaces. Additionally, we show that wider beams somewhat reduce median solution lengths (Fig. \ref{fig:bw_xgboost}B), indicating more optimal paths.

Finally, we assessed how computational resource demands scale with increasing permutation size when both beam width and non-backtracking depth are fixed at high-performing values. We found that both solution time and peak memory increase exponentially with $n$, revealing an asymptotic complexity of approximately $O(e^n)$ (Fig. \ref{fig:xgboost_time_mem}).

\begin{figure}[ht!]
  \centering
  \includegraphics[width=\textwidth]{figs/xgboost_time_mem_2panel.png}
  \caption{Solve time and peak memory as functions of permutation size $n$.
  Beam width and non-backtracking depth are fixed at high-performing values.
  The trends indicate approximately exponential growth with $n$ (complexity $\approx O(e^{n})$).}
  \label{fig:xgboost_time_mem}
\end{figure}

\clearpage
\subsection{Dependence on learning rate, epochs number and network size}
\label{sec-lr}
Here we present quite an interesting dependence on the combination of the  learning rate, epochs number and neural network size.
(More simulation results for various parameters can be found in \href{https://www.kaggle.com/competitions/lrx-oeis-a-186783-brainstorm-math-conjecture/discussion/612338}{spreadsheet here}, or  on \href{https://docs.google.com/spreadsheets/d/1KFYTqPDSdH8_vdvWz3YLTLKg5IlW1Kkx_9kJClSKGOs/edit?usp=sharing}{google sheets},
experiments are mostly done in the notebooks \href{https://www.kaggle.com/code/alexandervc/lrx-cayleypy-baseline2-mlp}{nptebook1}, \href{https://www.kaggle.com/code/alexandervc/lrx-cayleypy-rl-mdqn}{notebook2}.)
We mostly consider single layer MLP with layer size increasing by factor of 2 from 128 to 262144. 
Figures \ref{fig:SuccessRate1},\ref{fig:SuccessRate2} show success rates for $S_{16}$, beam size equal to 1,
pure diffusion distance training (without RL part). 
One can observe that, wider networks can achieve better results, however with smaller learning rate and larger number of epochs.
One also observe that optimal value of the learning rate is decreasing with number of epochs increased. 
Unfortunately we were unable to achieve 100\% success rate though results are quite close. 
Highlighted cells show best results in a row, blue highlight best result for table, yellow - if results are insignificantly different from the best, "lr" is learning rate. Number of experiments is typically 50-100, or shown in the brackets.
Typically one stops learning by when the validation loss stops to decrease, but in that case we observe that this standard strategy does not work, validation loss stabilizes quite early, however additional training improves the actual solution ability.
We attribute that to the noisy character of random walk data - loss cannot be made in zero in that situation in principle,
because each training batch (random walks) is different from the perfect diffusion distance. 
Thus we need to understand how many epochs to train for the given learning rate. 

\begin{figure}[ht!]
  \centering
  \includegraphics[width=\textwidth]{figs/SuccessRate1.png}
  \caption{Success Rate for 4000 and 16000 epochs.}
  \label{fig:SuccessRate1}
\end{figure}

\begin{figure}[ht!]
  \centering
  \includegraphics[width=\textwidth]{figs/SuccessRate2.png}
  \caption{Success Rate for 4000 and 16000 epochs.}
  \label{fig:SuccessRate2}
\end{figure}



\clearpage
\subsection{A Tropical Perspective on the Bellman Equation}\label{sec-tropical}
Tropical mathematics studies structures defined over a semiring equipped with the minimum and addition operations, instead of the usual addition and multiplication. The Bellman equation naturally embodies the tropical spirit due to the presence of the minimum over neighbors, in contrast to the averaging operations more common in graph theory. Here we discuss some elementary statements clarifying this connection.
%Generally speaking, there is an alternative approach to solving the classical Bellman equation. This leads us to \textit{tropical geometry}, a special field of mathematics whose methods are highly useful for studying algebraic geometry and graph theory. In this field, the problem of finding the shortest path in a weighted directed graph has a surprisingly simple solution. 

Consider a weighted directed graph $G$ with $n$ vertices. Its tropical adjacency matrix (\cite{MacSturm}, Chapter~1, Section~1.2) $A_{G},$ where the off-diagonal elements represent the lengths of the edges $a_{ij}\in\mathbf{R},$ and the diagonal elements are set to $a_{ii}=0.$ If there is no edge $(i, j)$ in $G,$ the corresponding entry in the matrix is assigned $a_{ij}=+\infty.$

{\bf Tropical $(n-1)$-power - lengths of shortest paths.}
It is easy to see (e.g. proposition~1.2.1 (\cite{MacSturm}, Chapter~1, Section~1.2) the tropical multiplication of the adjacency matrix $A_{G}$ with itself $(n - 1)$ times results in the matrix $P_{G}$ where the element $p_{ij}$ represents the lengths shortest path from the vertex $i$ to the vertex $j$ in $G.$ 

That is similar to the well-known fact, that matrix elements of standard $k$-powers of adjacency matrices give {\bf numbers} of paths of length $k$, but tropical powers give {\bf lengths} of paths, not their counting. 

Let $G$ be a connected graph with vertex set $\{v_{i}\}$ and edges of unit lengths.
%From the vertex form the set lets chose vertex $e_{k}.$ 
Then the $k$-th row of the matrix $P_{G}$ satisfies the Bellman equation (\ref{bel_eq}) with a boundary condition $d(v_{k}) = 0.$ In fact, the $(k,j)$ entry of this row equals the distance of the shortest path from $v_{k}$ to  $v_{j}$.

Notice that if the graph $G$ is connected, the matrix $P_{G}$ is represented by a weighted directed graph $G(P),$ that has a directed path from any vertex to any other vertex, which means that it is strongly connected. Then, Theorem~5.1.1 (\cite{MacSturm}, Chapter~5, Section~5.1) holds, stating that the matrix $P_{G}$ has only one eigenvalue, which is equal to 0.

Next, applying Theorem 5.1.3 (\cite{MacSturm}, Chapter~5, Section~5.1), we obtain that the columns of the matrix $P_{G}$ belong to the eigenspace of the tropical adjacency matrix $A_{G}.$ Thus, the solution to the Bellman equation can be associated with the problem of finding the eigenvectors of the matrix $A_{G}$ of the graph $G$ in tropical algebra.

Overall, one can think of the Bellman equation as a tropical analogue of the graph Laplacian, a central object of study in graph theory.
An intriguing question is: to what extent do the rich theories surrounding the Laplacian admit tropical counterparts?
In particular, since the cokernel of the Laplacian defines the sandpile (or Jacobian) group of a graph, one may ask whether there exists a tropical analogue of this construction.

\end{document} 